\documentclass[output=paper,
,modfonts
,nonflat]{langsci/langscibook} 

\ChapterDOI{10.5281/zenodo.1402549}

\title{The role of information structure for morphosyntactic choices in Tagalog}

\author{Anja Latrouite\affiliation{Heinrich-Heine-Universität Düsseldorf}\lastand  Arndt Riester\affiliation{Universität Stuttgart}}

% \chapterDOI{} %will be filled in at production
% \epigram{}

\abstract{In this paper we investigate the influence of two information structure (IS) related aspects on the choice of voice form and sentence structure by Tagalog speakers. The first is the information status of argument referents. Tagalog is a multiple voice language, so almost every semantic argument in a sentence can be turned into the privileged syntactic argument (or subject) and be rendered salient. Information status of the undergoer has been argued to play an important role in voice and subject selection. The second IS-related aspect is the inherent structure of a discourse as determined by the implicit questions under discussion (QUDs) that are answered with each subsequent sentence in a text. The default sentence in Tagalog starts with a verb. Inversion constructions, i.e.\ sentences that start with an argument phrase instead of a verb, are described as motivated by information structure considerations such as focus-background or contrastive-topic-focus packaging. Based on a novel QUD approach, we will work out the discourse structure and at-issue contents of five short texts and show the important role of implicit QUDs and parallelisms on the choice of voice and constituent order.}

\begin{document}
\maketitle

\section{Tagalog: Voice, information structure and inversion}\label{section_1}
\subsection{Multiple voice system}\label{section_voice}

\ili{Tagalog} is a verb-initial language known for its multiple \isi{voice system}. For an incremental theme verb like \textit{/sulat/} (\isi{actor voice}: \textit{sumulat}; \isi{undergoer voice}: \textit{sulatin}) ‘to write’ we find a set of \isi{voice} affixes on the verb that may pick out the semantic core arguments or one of the peripheral arguments as the privileged syntactic argument (PSA) -- the subject -- of the respective sentence. The PSA is marked by \textit{ang} if it is a common noun and \textit{si} if it is a personal name. The \textit{ang-} or \textit{si-}marked reference phrases (RPs, cf. \citealt{valro08}) tend to receive a definite or at least specific interpretation, and usually appear after the other core arguments, as shown in (\ref{write_basic}). Actor and \isi{undergoer voice} sentences differ in that undergoers\footnote{As it is common, we use the macro-role label \textit{undergoer} to refer to various non-actor roles.} expressed by common nouns preferably get an indefinite (and maybe even non-specific) reading in \isi{actor voice} sentences, e.g.\ \textit{ng liham} in (\ref{write_av1}), while they preferably receive a definite or specific reading in \isi{undergoer voice} sentences, e.g.\ \textit{ang liham} in (\ref{write_tv1}).\footnote{The (specific) indefinite reading of the \isi{undergoer} can be enforced by inserting the numeral \textit{isa-ng} `one'.

\begin{exe}
\ex \textbf{–in: (incremental) theme voice}\\
\gll Hindi s{\USSmaller}in{\USGreater}ulat ng babae kay Lisa \textbf{ang} \textbf{isa-ng} \textbf{liham}.\\
 	\textsc{neg}  \textsc{<uv.rls>}write \textsc{gen} woman \textsc{dat} Lisa \textsc{nom} one-\textsc{lk} letter\\
\glt `The woman did not write \textbf{a certain letter} to Lisa.'
\end{exe}
} Note that the actor tends to receive a definite reading regardless of case marking.

\begin{exe}
\ex\label{write_basic} /sulat/ ‘write’
\begin{xlist}
\ex\label{write_av1} \textbf{um-: actor voice}\\
\gll Hindi s{\USSmaller}\textbf{um}{\USGreater}ulat {{\ob}ng} {liham{\cb}\ug} kay Lisa \textbf{{\ob}ang} \textbf{babae{\cb}\act}.\\
     \textsc{neg}   \textsc{<av.rls>}write \textsc{\void{[}gen} letter  \textsc{dat} Lisa \textsc{\void{[}nom} woman\\ 
\glt `\textbf{The woman} did not write a(ny) letter/letters to Lisa.'
\ex\label{write_tv1} \textbf{–in: (incremental) theme voice}\\
\gll Hindi s{\USSmaller}\textbf{in}{\USGreater}ulat {{\ob}ng} {babae{\cb}\act} kay Lisa \textbf{{\ob}ang} \textbf{liham{\cb}\ug}.\\
 	\textsc{neg}  \textsc{<uv.rls>}write \textsc{\void{[}gen} woman \textsc{dat} Lisa \textsc{\void{[}nom} letter\\
\glt `The woman did not write \textbf{the letter/a (certain) letter} to Lisa.'
\end{xlist}
\end{exe}

\noindent As mentioned above, it is not only the agent or the theme argument that may be singled out as salient with a verb like \textit{write}, as the examples in (\ref{peripheral1}) show. Note that if an argument other than the agent or the theme is chosen as PSA, the \isi{information status} of the theme argument is unspecified. 
Sentence (\ref{write_ttv1}) exemplifies that thematic role labels (here: \textit{theme}) are sometimes not fine-grained enough to capture the semantic properties that the \isi{voice} affixes pick out. We will not be concerned with the intricacies of so-called peripheral \isi{voice} forms in this paper. The examples are only meant to show that the \isi{voice system} serves first and foremost to make different semantic arguments in a sentence salient and that only one argument at a time may be marked as salient in a sentence. 

\begin{exe}
\ex\label{peripheral1}
\begin{xlist}
\ex\label{write_gv1} \textbf{–an: goal voice}\\
\gll S{\USSmaller}in{\USGreater}ulat-\textbf{an}    ng babae    ng liham    \textbf{si} \textbf{Lisa}.\\
     <\textsc{uv.rls}>write‐\textsc{gv} \textsc{gen} woman   \textsc{gen} letter \textsc{nom} Lisa\\ 
\glt `The woman wrote \textbf{Lisa} a/the letter/letters.'
\ex\label{write_ttv1} \textbf{i-: (transferred) theme voice}\\
\gll \textbf{I}‐s{\USSmaller}in{\USGreater}ulat    ng babae    \textbf{ang} \textbf{pangalan}    \textbf{ni} \textbf{Lisa} sa papel.\\
 \textsc{tv-<uv.rls>}write \textsc{gen} woman \textsc{nom} name   \textsc{gen} Lisa  \textsc{dat} paper\\
\glt `The woman wrote \textbf{Lisa’s name} on (the) paper.'
\ex\label{write_iv1} \textbf{i-pan(g)-: instrument voice}\\
\gll \textbf{I-p}{\USSmaller}in{\USGreater}\textbf{an}‐sulat         ng babae     kay Lisa    ng liham    \textbf{ang} \textbf{lapis} \textbf{mo}.\\ 
	\textsc{iv-appl$_\text{instr}$<uv.rls>}-write \textsc{gen} woman \textsc{dat} Lisa  \textsc{gen} letter \textsc{nom} pencil your\\
\glt `The woman wrote the/a letter/letters to Lisa \textbf{with your pencil}.'
\end{xlist}
\end{exe}

\noindent So far there are no elaborate corpus studies exploring whether the de\-fi\-nite\-ness/\-spe\-ci\-fi\-ci\-ty associated with \textit{ang}-marking is about uniqueness, identifiability, familiarity, \isi{givenness}, a certain kind of accessibility or other in\-for\-ma\-tion-struc\-tu\-ral\-ly relevant features. In this paper we are going to approach the problem from the other side. We will analyze the information structure and \isi{discourse structure} of textual data and subsequently study the influence of information-structural constellations on \isi{voice selection}. By \textit{constellation} we mean, in particular, the focus-background structure of the entire \isi{clause} within its \isi{discourse} context rather than isolated properties such as the \isi{givenness} of a single argument. The goal is to make further contributions to the investigation of which \isi{discourse} properties exactly motivate \isi{voice} and construction selection in a given text. 
To this end we will look at the results of a \ili{Tagalog} translation and evaluation study aimed at getting speakers to produce and rate the same sentence in five contexts that differ with respect to the \isi{givenness} of the arguments and the overall \isi{discourse structure}. 
Before we turn to this study, however, a few more words need to be said about what has been found so far with respect to \isi{voice}, \isi{constituent order} and  information structure coding in \ili{Tagalog}. 

\subsection{Preliminary assumptions on information status, event type, information structure and the prominence of the undergoer}\label{sect:pat_prom}

Philippine languages have been characterized as patient-prominent languages by \cite{cenre77}, \cite{guzvi92} and others, in the sense that \isi{undergoer voice} has been said to be more frequent than \isi{actor voice} in text counts, at least in transitive sentences \citep[cf.][]{payto94,selpe01}. \cite{macla02} among others argues that this observation does not hold unequivocally for \ili{Tagalog}. Investigating several Philippine languages, \cite{nolri05}, who proposes an ergative analysis of \ili{Tagalog}, characterizes Philippine ergativity as speakers giving "the highest degree of \isi{prominence} to the most affected entity" (ibid, p.\ 236), i.e.\ the entity viewed as most saliently affected by the event is said to receive the absolutive\footnote{Since we do not subscribe to the ergative analysis, we will gloss \textit{ang} as \textit{nominative} instead.} case (\textit{ang-}marking). For the most part, the most affected argument in transitive scenarios is the \isi{undergoer}, but may be the actor if the \isi{undergoer} is less identifiable. So, basically, he also recurs to the degree of \isi{referentiality} (i.e.\ the \isi{information status}) of the \isi{undergoer} as one of the essential factors for \isi{voice selection}. The idea of \textit{ang-}marking as \isi{prominence} marking is taken up by \cite{latan11}, who suggests three levels of evaluation with respect to the relative \isi{prominence} of an argument: (i) the referential level, (ii) the event-structural level, and (iii) the information-structural level. 

With respect to the referential level there are two possible scenarios: (i) one argument is higher on the scale of \isi{referentiality} than the other (given > familiar > unfamiliar > non-unique > non-specific)\footnote{Compare also the \isi{givenness} hierarchy by \cite{gunje93} and the comprehensive overview on \isi{information status} by \cite{baust12a}.} and, therefore, considered to be more prominent, or (ii) both are equally referential, but the \isi{referentiality} of one of the arguments is less expected and, therefore, this information is treated as more salient. Expectancy is based on the frequency of a pattern. As \cite{pribe12} and others have pointed out, actors tend to be referentially independent and definite, while undergoers tend to be referentially dependent on the actor and are less often definite. Therefore, we find cross-linguistically a tendency to develop marked morphosyntactic coding for definite undergoers.

With respect to the level of event-structural \isi{prominence}, we can distinguish un\-der\-goer- from actor-prominent and neutral verbs. Clearly undergoer-oriented verbs are change-of-state verbs like \textit{to scare} or \textit{to kill} that only provide information on the result state of the \isi{undergoer}, but no information regarding the specific activity of the actor. Actor-oriented verbs are manner of action verbs (e.g.\ \textit{to devour}) and certain stative verbs (e.g.\ \textit{to fear}), and neutral verbs comprise punctual contact verbs like  \textit{to hit}. \cite{schpa72} provide examples showing that an undergoer-oriented verb like \textit{to kill} or \textit{to scare} cannot be used with \isi{actor voice} in an unmarked, verb-initial sentence, as exemplified in (\ref{scare1}), but only with \isi{undergoer voice}, as in (\ref{scare1b}). 

\begin{exe}
\ex
\begin{xlist}
\ex[*]{\label{scare1}
\gll T{\USSmaller}um{\USGreater}akot   sa    mga babae si Lena.\\
     <\textsc{av.rls}>scare    \textsc{dat} \textsc{pl}     woman \textsc{nom}  Lena\\
\glt  Intended: `Lena scared the women.'}
\ex[\hphantom{*}]{\label{scare1b}
\gll T{\USSmaller}in{\USGreater}akot    ni Lena   ang isa-ng babae.\\
     <\textsc{uv.rls}>scare    \textsc{gen} Lena \textsc{nom}  \textsc{one-lk}  woman\\
\glt`Lena scared a woman.'}
\end{xlist}
\end{exe}

\noindent This suggests that the \isi{voice} alternation we are interested in is not available for these verbs. Note, however, that there are examples of \isi{actor voice} forms of these verbs, even with definite undergoers, as example (\ref{scare2}) shows. The sequence in (\ref{scareEx2}) is taken from a blog discussing a movie. The movie is about sisters that suffer from different kinds of traumas. 

\begin{exe}
\ex\label{scareEx2}
\begin{xlist}
\ex\label{develop1}
\gll Kay Angel, ang multo ay ang pagiging      alipin sa alak at sigarilyo.\\
 \textsc{dat} Angel, \textsc{nom} ghost   \textsc{inv} \textsc{nom} developing slave   \textsc{dat} alcohol and cigarette\\
\glt `For Angel, the ghost is the developing addiction to alcohol and cigarettes.'
\ex\label{scare2}
\gll {{\ob}Kay} {Sylvia,{\cb}\ug} t{\USSmaller}\textbf{um}{\USGreater}a-takot {{\ob}sa} {kaniya{\cb}\ug} {{\ob}ang} pagpapalaglag sa kanya-ng magiging {anak{\cb}\act}.\\
  	\textsc{\void{[}dat} Sylvia \textsc{<av.rls>ipfv}.scare \textsc{\void{[}dat} her \textsc{\void{[}nom} abortion \textsc{dat} her-\textsc{lk} developing child\\
\glt `As for Sylvia, the abortion of her developing child scares (is haunting) her.'\\
\hfill \texttt{http://panitikanngpilipinas17.blogspot.de}
\end{xlist}
\end{exe}

\noindent \cite{latan11} observes that some kind of event-related markedness seems to be necessary to license \isi{actor voice} of an inherently undergoer-oriented verb in a verb-initial \isi{clause}. In her examples, the verb is always marked for imperfective aspect, inducing a habitual reading, and the \isi{actor argument} is an inanimate causer rather than a volitional, controlling agent. She suggests that it is these divergent properties that license the use of the marked \isi{actor voice} form. Note that in the example given in (\ref{scare2}), the sentence is also in\-for\-ma\-tion-struc\-tu\-ral\-ly marked, as the \isi{undergoer} is a \isi{contrastive topic} and appears in the \isi{left periphery}. The new, focal  information in (\ref{scare2}) is, thus, provided by the actor phrase \textit{the abortion of her developing child}. The example is couched in a series of contrasting sentences, one of which is given in (\ref{develop1}), listing the afflictions and problems the -- previously introduced -- protagonists are haunted by. \cite{latan11} suggests that considerations with respect to the level of in\-for\-ma\-tion-struc\-tu\-ral pro\-mi\-nence outrank considerations regarding the other two levels (event type and \isi{information status} of arguments). Based on examples like (\ref{scare2}) as well as examples discussed further down below, it is argued that information structure plays an important role with respect to \isi{voice selection}. However, there is no systematic corpus study to back-up the proposal as of yet and no detailed in\-for\-ma\-tion-struc\-tur\-al analysis of the data. This paper is meant to further investigate the claim that not only \isi{givenness} and \isi{topicality}, but also focality plays an important role for \isi{voice selection} and, therefore, for PSA-marking. 

Note that for the two not very frequent peripheral \isi{voice} forms, instrument \isi{voice} and causative \isi{voice}, \cite{nuhpa16} finds textual \isi{givenness} of the PSA to be a prerequisite in his preliminary corpus study. It seems to make sense that \isi{undergoer} arguments can only be realized as the PSA if they are prementioned. However, note that the licensing of \isi{actor voice} in (\ref{scare2}) rather points to focality of the actor as the decisive factor, since it is chosen over a given \isi{undergoer}.

\cite{latan16a} finds that the textual \isi{givenness} of undergoers is not sufficient to make speakers choose \isi{undergoer voice}. For instance, in the target sentence (\ref{given_u_target1}) the prementioned \isi{undergoer} (experiencer) \textit{mga negosyante} `the negotiators' (see context in \ref{given_u_context1}) is not marked by \textit{ang}; rather it is the inanimate \isi{actor argument} \textit{the rallies} which receives \textit{ang}. As we have just pointed out, \textit{/takot/ (tumakot, takutin)} `to scare' is clearly undergoer-oriented. So not only based on the \isi{givenness} of the \isi{undergoer} but also based on event-structural \isi{prominence} considerations, the \isi{undergoer} argument would be expected to turn into the PSA. However, once again, we find \isi{actor voice} and, once again, the sentence is information-structurally marked. It negates the truth of the previous sentence, and the follow-up sentence shows that the reason for the negotiators’ fear is not the rallies but Erap.\footnote{\textit{Erap} is the nickname of Joseph Estrada, at the time President of the Philippines. Note, furthermore, that the form \textit{yung} is very often used instead of \textit{ang} in spoken \ili{Tagalog} and seems to find its way into written \ili{Tagalog} as well.} It is therefore possible that it is the contrast between the actor arguments that licenses the chosen \isi{actor voice} form.

\newpage 
\begin{exe}
\ex\label{given_undergoer1} Actor \isi{voice} despite prementioned \isi{undergoer}:\\
\begin{xlist}
\ex\label{given_u_context1}\textbf{Context (\isi{undergoer voice}) providing \textit{rallies} and \textit{negosyante}:}\\
\gll Sabat naman ni Executive Sec. Ronaldo Zamora, t{\USSmaller}in{\USGreater}a-takot lang ng rallies ang mga negosyante.\\
answer likewise \textsc{gen} Executive Sec. Ronaldo Zamora \textsc{<uv>ipfv}-scare    just    \textsc{gen} rallies \textsc{nom} \textsc{pl} negotiators\\
\glt `According to Executive Sec.\ Ronaldo Zamora, the rallies just frighten the negotiators/businessmen.' 
\ex{\textbf{Additional context:}}\\
`And according to the postscript by Education Sec.\ Andrew Gonzales, it is forbidden for teachers to accompany the students to the rally, and a discussion regarding impeachment (from class) comes on top [\dots] What is this, martial law?'
\ex\label{given_u_target1}\textbf{Target sentence:}\\
\gll Hindi totoo-ng t{\USSmaller}\textbf{um}{\USGreater}a-takot {{\ob}ng} mga {negosyante{\cb}\ug} {{\ob}ang} {rallies{\cb}\act}.\\
Not true-\textsc{lk} \textsc{<av.rls>ipfv}-scare \textsc{\void{[}gen} \textsc{pl} negotiator \textsc{\void{[}nom} rallies\\
\glt `It is not true that the rallies frighten the negotiators/businessmen.'
\ex\label{given_u_cont1}\textbf{Continuation:}\\
\gll Ang t{\USSmaller}\textbf{um}{\USGreater}a-takot {{\ob}sa} {kanila{\cb}\ug,} {{\ob}yung} pananatili ni Erap sa {puwesto{\cb}\act}.\\
  \textsc{nom} \textsc{<av.rls>ipfv}-scare \textsc{\void{[}dat} them \textsc{\void{[}nom} staying \textsc{gen} Erap \textsc{dat} spot\\
\glt `What scares them is Erap’s staying on his spot.'\\
\hfill (Pilipino Star Ngayon, December 12, 2000, \textit{Mag -- rally or tumahimik})
\end{xlist}
\end{exe}

\noindent The sentence in (\ref{hunger_games1}) shows an example pointing in a similar direction. The example is taken from the \ili{Tagalog} translation of the \textit{Hunger Games} by Suzanne \citet{collins09}. In the preceding paragraph the reader learns that the narrator is on her way to the woods to go hunting, which is illegal and may result in her getting killed. The narrator muses about the dangers she is putting herself into and how she cannot be open and truthful with anyone, not even her closest family members. In this context the following sentence is uttered.\footnote{\ili{English} original: "In the woods waits the only person with whom I can be myself: Gale."}

\begin{exe}
\ex\label{hunger_games1}
\gll Nag-hi-hintay {{\ob}sa} {akin{\cb}\ug} sa gubat {{\ob}ang} nagiisa-ng tao-ng puwede ko-ng pa-kita-an ng tunay na ako -- si {Gale{\cb}\act}.\\
\textsc{av.rls-ipfv}-wait \textsc{\void{[}dat} \textsc{1sg} \textsc{dat} forest \textsc{\void{[}nom} alone-\textsc{lk} man-\textsc{lk} can   \textsc{1sg.gen-lk}    \textsc{caus}-visible-\textsc{lv}    \textsc{gen} true \textsc{lk} \textsc{1sg.nom} {}    \textsc{nom} Gale\\
\glt `Waiting for me in the woods is the only person I can show
my true self to -- Gale'  (\textit{Yung Hunger Games}, p.\ 4)
\end{exe}

\noindent The actor phrase \textit{the only person [I can show my true self to]: Gale} is newly introduced and, based on the \isi{relative clause}, also starkly contrasted to the previously mentioned family members that the narrator cannot confide in. The \isi{undergoer} argument is the narrator herself, so she is clearly given. Despite this fact, \isi{actor voice} is chosen and the new, contrasted \isi{actor argument} Gale is chosen as subject.\footnote{Note that without further context the \isi{undergoer voice} form would be preferred by speakers if the \isi{undergoer} is highly referential (active, salient, given).
\begin{exe}
\ex
\gll H{\USSmaller}in{\USGreater}intay ako ng  isa-ng lalaki {\USSlash} ni Gale. \\
     \textsc{<uv.rls>}wait \textsc{1sg.nom}  \textsc{gen} one-\textsc{lk} man {} \textsc{gen} Gale\\
\glt `A man/Gale waited for me.'
\end{exe}} Examples like these lead \cite{latan11} to the conclusion that the focality of the \isi{actor argument} may result in actor \isi{voice selection}, a pattern that cannot be predicted from the \isi{information status} of the \isi{undergoer} alone. It has been noted before that there is a \isi{default mapping} between \isi{givenness}/\isi{topicality} and the macro-role \textit{actor} on the one hand, and newness/focality and the macro-role \textit{undergoer} on the other hand (e.g. \citealt{Lambrecht1994} among others), and that divergence from this \isi{default mapping} often results in a specific morphosyntactic marking in the languages of the world, cf.\ \cite{gueto15} on \ili{African} languages.\footnote{The reason for this \isi{default mapping}, at least in narratives, may very well be that stories tend to center around a small number of protagonists that engage in various activities with objects and people they encounter in the course of a story, so that the \isi{information packaging} of a default sentence is Actor (topical) Verb (focal) \isi{Undergoer} (focal).} Therefore, it is worthwhile exploring whether the sentences above are exceptions or whether they point to a more systematic pattern.

\subsection{Hypotheses on information structure and voice marking}\label{sect:hypo}

Given the asymmetry with respect to the \isi{default mapping} of macro-roles and in\-for\-ma\-tion-struc\-tu\-ral values, the characterization of in\-for\-ma\-tion-struc\-tu\-ral pro\-mi\-nence needs to be stated as follows:

\begin{exe}
\ex\label{is_prom} \textbf{Information-structural \isi{prominence} (characterization):}\\
A core argument is IS-prominent if it has a non-default IS value (whereby the default value for actors is \textit{+topical} and the default value for undergoers is \textit{+focal} in a transitive sentence).
\end{exe}

\noindent Based on the \isi{actor focus} examples above and cross-linguistic findings that in\-for\-ma\-tion-structural \isi{prominence} as characterized in (\ref{is_prom}) is often-times reflected morphosyntactically in the languages of the world, we can formulate the hypothesis that in\-for\-ma\-tion-struc\-tu\-ral \isi{prominence} may indeed be a decisive factor for \isi{voice selection} in \ili{Tagalog}.  If \isi{voice selection} is influenced by \isi{prominence} considerations and IS-\isi{prominence} is an important factor, we expect the following correlation:

\begin{exe}
\ex\label{tendencies} IS-influenced tendencies w.r.t.\ \isi{voice selection} with two-place predicates:
\begin{xlist}
\ex If the \isi{undergoer} is topical (given), \isi{undergoer voice} is preferred.
\ex If the actor is focal (new), \isi{actor voice} is preferred
\end{xlist}
\end{exe}

\noindent It is important to note that research on information structure of the past decades has made it reasonably clear that the conceptual pairs \textit{focal} vs.\ \textit{new} (as well as \textit{topical} vs.\ \textit{given}) are closely related  yet not identical to each other, see e.g.\ \cite{beada11} or \cite{riear13}. This is why, in \sectref{sect:is_theory}, we are going to switch to a question-based definition of focus. In the current section, however, and with regard to the examples introduced so far, a novelty-based definition of focus is sufficient and, perhaps, easier to comprehend.

The tendencies stated in (\ref{tendencies}) render clear predictions if the actor and the \isi{undergoer} do not differ in status with respect to information structure, i.e.\ if both are topical or both are focal. If both are topical, only the \isi{undergoer} is considered information-structurally prominent based on (\ref{is_prom}). If both are focal, only the actor is considered %information-structurally 
prominent. As a first hypothesis to be checked we suggest the following:

\hypo{i}{If the actor is focal (contrary to the \isi{default mapping}) and the \isi{undergoer} is focal (in accordance with the \isi{default mapping}), then \isi{actor voice} is chosen:\\[1ex] Actor\focus\  \isi{Undergoer}\focus\ $\rightarrow$ Actor Voice\hfill(preliminary)
}

\noindent The situation described in Hypothesis (i) that will come to mind at first is probably so-called \textit{broad focus} although the hypothesis can, for instance, also be applied to situations in which the \isi{undergoer} is a (focus-like) \isi{contrastive topic} and the actor is a focus, as witnessed in example (\ref{scare2}). We will say more on \isi{contrastive} topics in \sectref{sect:is_tagalog}. A second hypothesis to be checked is:

\hypo{ii}{If the actor is topical (according to the \isi{default mapping}) and the \isi{undergoer} is topical (contrary to the \isi{default mapping}), then \isi{undergoer voice} is chosen:\\[1ex]
Actor\topic\  \isi{Undergoer}\topic\ $\rightarrow$ \isi{Undergoer} Voice}

\noindent \isi{Undergoer} \isi{voice} could then be expected to be the preferred choice in sentences in which the verb or some non-core argument is focal, but both actor and \isi{undergoer} are old information.

There are two more possible scenarios. The first scenario is the \isi{default mapping}: actor (topical) -- \isi{undergoer} (focal). Neither argument is in\-for\-ma\-tion-struc\-tu\-ral\-ly prominent according to (\ref{is_prom}) in this scenario. Therefore, the choice of the respective \isi{voice} form will have to be based on \isi{prominence} considerations at a different level. The default for discourse-new undergoers is that they are indefinite, i.e.\ not referentially prominent. Based on the referential \isi{prominence} considerations mentioned above, we expect \isi{actor voice} to be frequent in this scenario, as shown, for instance, in (\ref{write_av1}); cf.\ \cite{pribe12} on definiteness as a generally marked option for undergoers/objects. However, \isi{undergoer voice} is also found in this constellation; compare (\ref{scare1b}). Therefore, the only claim made in Hypothesis (iii) is that \isi{voice selection} in this scenario involves information from a different level.

\newpage 
\hypo{iii}{If the actor is topical (according to the \isi{default mapping}) and the \isi{undergoer} is focal (according to the \isi{default mapping}), then both \isi{voice} forms may be used. The final choice will depend on \isi{prominence} considerations at a different level (i.e.\ referential \isi{prominence} or event-structural \isi{prominence}).\\[1ex]
Actor\topic\ \isi{Undergoer}\focus\ $\rightarrow$ Actor Voice $\vee$ \isi{Undergoer} Voice}

\noindent The last possible scenario is a very marked one: both actor and \isi{undergoer} are in\-for\-ma\-tion-struc\-tur\-al\-ly prominent given the characterization in (\ref{is_prom}). In view of the fact that a given \isi{undergoer} is not only IS-prominent but also prominent at the referential level, we might expect \isi{undergoer voice} to be preferred over \isi{actor voice} in this case. However, as we have already seen in examples (\ref{given_u_cont1}) and (\ref{hunger_games1}) above, a given \isi{undergoer} does not necessarily enforce \isi{undergoer voice}, so focality of the actor seems to outrank the \isi{topicality} of the \isi{undergoer} in quite a few instances. 

\hypo{iv}{Actor\focus\ %(+C)
\isi{Undergoer}\topic\ $\rightarrow$ Actor \isi{voice} > \isi{Undergoer} Voice\hfill(preliminary)}

\noindent Morphological marking on the verb is not the only possibility to overtly express IS-\isi{prominence}. Many languages use fronting and inversion constructions as well as \isi{prosody} to mark IS-\isi{prominence}, and so does \ili{Tagalog}. Therefore, we can expect that speakers may be able to express the IS-\isi{prominence} of both arguments, albeit with different means: for instance, \isi{voice} marking for one of the arguments and inversion for the other one. This is indeed found, as discussed in \sectref{sect:is_tagalog}.

In order to factor out event-structural \isi{prominence} and to highlight the role of information structure, \cite{latan16a} construed contexts for one and the same target sentence, containing the same verb and reference phrases. The contexts were meant to clearly determine the \isi{information status} of the reference phrases and narrow down the set of possible implicit questions the target sentence could be an answer to. In this chapter, these contexts will be examined from a discourse-structural perspective. Before we turn to the study, however, we need to take a look at \isi{constituent order} and information structure in \ili{Tagalog} and lay out our question-under-discussion approach which we utilize to determine the information structure of an utterance.

\subsection{Constituent order, prosody and inversion constructions}\label{sect:is_tagalog}

In this section, we turn our attention to other means of information-structure marking than \isi{voice}. \cite{kauda05} puts forward what he calls the ``double focus" construction in (\ref{kaufman1b}), in which only the (non-focal) \isi{actor argument}\footnote{ (\ref{kaufman1b}) shows that a pronominal (in contrast to a nominal) PSA does not occur sentence-finally and is not \textit{ang}-marked.} is signalled on the verb via \isi{voice} morphology, while the two focal \isi{undergoer} arguments in-situ are marked by prosodic stress. Given what has been just laid out in \sectref{sect:hypo}, we would not expect focal undergoers to trigger \isi{undergoer voice}, so the data do not yet pose a challenge. There have been no claims so far with respect to how contrast should be marked. Prosody (indicated by capital letters) seems to be a possible option. 

\begin{exe}
\ex\label{kaufmannex1} \textbf{``Double focus" \citep[187ff.]{kauda05}} 
\begin{xlist}
\ex
\gll Dapat    ba-ng    mag-bigay    ng pera    sa mga guro{\USQMark}\\
            should    \textsc{q-lk}    \textsc{av}-give    \textsc{gen} money    \textsc{dat} \textsc{pl} teacher\\
\glt `Should one give money to the teachers?'
\ex \label{kaufman1b}
\gll Hinde, mag-bigay ka {na lang} {{\ob}ng} {KENDI{\cb}\focus} {{\ob}sa} mga {BATA{\cb}\focus}.\\
\textsc{neg} \textsc{av}-give \textsc{2sg.nom} only \textsc{\void{[}gen} candy \textsc{\void{[}dat} \textsc{pl} child\\
\glt `No, (you) just give CANDY to the KIDS.' 
\end{xlist}
\end{exe}

\noindent Very often the phenomenon that some people call ``double focus" is actually a combination of a \isi{contrastive topic} \citep{bueda03} plus a (\isi{contrastive}) focus.\footnote{True instances of ``double focus", which may also be characterized as two parts of a single, discontinuous \textit{complex focus} \citep{krima92}, are rare in comparison with the rather frequent \textsc{ct-f} pairs. See also the discussion in \cite[216]{riear13}.} Contrastive topics have pragmatically a lot in common with foci (both of them give rise to alternatives, cf. \citealt{rooma92}). We therefore expect that our hypotheses made with respect to focality (in particular, Hypothesis i) also cover \isi{contrastive} topics. For instance, in (\ref{scare2}) we already saw actor-\isi{voice} marking in combination with a \isi{contrastive topic} and a focus. 

A rather special means of focus/contrast marking is shown in example (\ref{martinez_b}), by \cite[194]{kauda05}: the \isi{contrastive}, genitive-marked actor appears in an extraordinary position at the end of the sentence and receives prosodic stress, while the \isi{undergoer} is given. Thus, according to Hypothesis (iv), we would expect \isi{actor voice}. However, in (\ref{martinez_b}) \isi{undergoer voice} is found, indicating that special \isi{word order} choices seem to have an influence on \isi{voice selection} as well.

\begin{exe}
\ex
\begin{xlist}
\ex\label{martinez_a}
\gll K{\USSmaller}in{\USGreater}a-usap ng bawat propeso ang mga estudyante, di ba{\USQMark}\\   
     \textsc{<uv>com}-speak \textsc{gen} each professor \textsc{nom} \textsc{pl} student, \textsc{neg} \textsc{q}\\
\glt `Each professor spoke with the students, right?'
\ex\label{martinez_b}
\gll Hindi. Hindi k{\USSmaller}in{\USGreater}a-usap ang mga estudyante ni {{\ob}PROPESOR} {MARTINEZ{\cb}\focus}.\\
     \textsc{neg} \textsc{neg} \textsc{<uv>com}-speak \textsc{nom} \textsc{pl}     student \textsc{gen}  {\void{[}Prof.} Martinez\\ 
\glt `No. PROFESSOR MARTINEZ did NOT speak with the students.' 
\end{xlist}
\end{exe}

\noindent Note that the actual information structure of (\ref{martinez_b}) is in fact a bit unclear. It is quite likely that, other than what is assumed in \cite{kauda05}, the extraposed phrase \textit{ni Propesor Martinez} is actually the \isi{contrastive topic}, while the focus is on the negation. Since we have no other piece of evidence for this particular construction, we shall have nothing more to say about it at this point.

It is also possible for a contrasted core argument to appear with \isi{nominative} marking in sentence-initial position followed by \textit{ay}, while the other focal element appears in-situ with prosodic stress, as shown in (\ref{ct_f1}). This construction is one of three inversion constructions mentioned in \cite{schpa72} and \cite{valro15} and has also been linked to \isi{information packaging} considerations. Inversion constructions are structures that diverge from the verb-initial default \isi{word order} by starting off with a reference phrase in sentence‐initial position. The three kinds of inversion constructions are the so-called \textit{ay}-inversion construction, the \textit{ang}-inversion construction and the adjunct fronting construction. In our study below, the \textit{ay}-inversion construction was the most frequent one. According to \cite{schpa72}, \textit{ay}-inverted elements are for the most part topical arguments, but may be used for contrast. In the example in (\ref{ct_f1}), the \textit{ay}-inverted element can easily be construed as a \isi{contrastive topic} and the argument realized in-situ as the (\isi{contrastive}) focus. 

\begin{exe}
\ex\label{ct_f1} \textbf{Contrastive topic -- \isi{contrastive} focus}\\
\gll {{\ob}Si} {Peter{\cb}\ctopic} ay nag-ba-basa {{\ob}ng} LIBRO {mo{\cb}\focus.}\\
\textsc{\void{[}nom} Peter \textsc{inv} \textsc{av.rls-ipfv}-read    \textsc{\void{[}gen} book  \textsc{2sg.gen}\\
\glt ‘PETER is reading your BOOK.’
\end{exe}

\noindent Note that, despite its \isi{nominative} marking, the sentence-initial element followed by \textit{ay}  does not have to be the PSA of the verb; at least in \isi{undergoer voice} sentences it can be the actor as well, see (\ref{ay_inv_1}). The particle \textit{ay} may also be replaced by a pause here, signalling the clause-external position of the sentence-initial element. There are no corpus studies yet, but one can speculate that the construction in (\ref{ay_inv_1}) may be one of the ways to code both the \isi{topicality} of the \isi{undergoer} (\isi{undergoer voice}) and the focality of the actor (\textit{ay}-inversion) at the same time. We therefore tentatively assume the information structure indicated on the example. Note that it is not possible to \textit{ay}-invert an \isi{undergoer} in combination with \isi{actor voice}, see (\ref{ay_inv_2}).\footnote{Even with a resumptive \isi{pronoun}, the sentence is considered ungrammatical by the consultants.

\begin{exe}
\ex[*]{
\gll Ang libro mo   ay    nag-ba-basa  nito si       Peter.\\
     \textsc{nom} book your \textsc{inv} \textsc{<av.rls>}read this.\textsc{gen} \textsc{nom} Peter\\
\glt Intended: `The book, Peter is reading it.'}
\end{exe}
} 

\begin{exe}
\ex
\begin{xlist}
\ex[\hphantom{*}]{\label{ay_inv_1} \textbf{\textit{Ay}-inversion sentence with two \textsc{nom}-marked phrases}\\
\gll {{\ob}Si} {Peter{\cb}\focus} ay b{\USSmaller}in{\USGreater}asa {{\ob}ang} libro {mo{\cb}\topic}.\\
\textsc{\void{[}nom} Peter \textsc{inv} \textsc{<uv.rls>}read   \textsc{\void{[}nom} book your\\
\glt `PETER read your book.'}
\ex[*]{\label{ay_inv_2}
\gll Ang libro mo   ay    nag-basa   si       Peter.\\
     \textsc{nom} book your \textsc{inv} \textsc{<av.rls>}read \textsc{nom} Peter\\
\glt Intended: `Peter read the BOOK.'}
\end{xlist}
\end{exe}

\noindent Based on these descriptions we can add a fifth hypothesis that we wish to check.

\hypo{v}{Actor\focus\ +ay \isi{Undergoer}\topic\ $\rightarrow$ \isi{Undergoer} Voice\hfill(preliminary)} 

\noindent Meanwhile, based on the problems with extraposition in (\ref{martinez_b}) and with \textit{ay}-in\-ver\-sion in (\ref{ay_inv_1}), we change Hypothesis (iv) to exclude \textit{ay}-inversion and non-default \isi{constituent order}.

\hypo{iv}{\isi{Undergoer}\topic\ Actor\focus\ \textit{(only default \isi{constituent order})} $\rightarrow$ Actor Voice\hfill(final)}

\noindent Finally, Hypothesis (i) will be stated more precisely to explicitly also cover \textit{ay}-inversion and \isi{contrastive} topics, i.e.\ examples like (\ref{ct_f1}).

\hypo{i}{Actor$_{\text{\footnotesize F/CT}}$ \wupperhand ay \isi{Undergoer}\focus\ $\rightarrow$ Actor Voice\hfill(final)} 

\noindent In the following, we mention a few more information-structurally relevant syntactic properties of \ili{Tagalog}, which, however, we will not investigate any further. It should be noted that the \textit{ay}-inversion may also be used in connection with framesetters, such as a temporal or local adverbial phrase, as shown in (\ref{framesetter}).

\begin{exe}
\ex\label{framesetter}
\gll Kahapon    ay    nag‐basa    ng libro mo    si Peter.\\
     yesterday    \textsc{inv}    \textsc{av.rls}‐read    \textsc{gen} book your  \textsc{nom} Peter\\
\glt `Yesterday Peter read your book.'
\end{exe}

\noindent If \ili{Tagalog} speakers wish to put \isi{narrow focus} on an obliquely marked argument or adjunct, they are also said to recur to adjunct inversion \citep{kropa93}, in which the adjunct is moved to the sentence‐initial position and attracts clitics (if present) that appear between them and the verb. In (\ref{adj_inv_1}) the \isi{pronoun} \textit{siya} is a \isi{clitic}.

\begin{exe}
\ex\label{adj_inv_1} \textbf{Adjunct/oblique inversion (``emphatic" inversion, Schachter \& Otanes 1972, pp.\ 496--98)}
\begin{xlist}
\ex
\gll T{\USSmaller}um{\USGreater}awa    siya    sa kaniya      kahapon.\\
     <\textsc{av.rls}>laugh   \textsc{3sg.nom}  \textsc{dat} \textsc{3sg.obl} yesterday\\
\glt `She laughed at him yesterday.'
\ex
\gll Kahapon    siya    t{\USSmaller}um{\USGreater}awa    sa kaniya.\\   
yesterday    \textsc{3sg.nom}    <\textsc{av.rls}>laugh   \textsc{dat} \textsc{3sg.obl}\\  
\glt `It was yesterday that she laughed at him.'
\ex
\gll Sa kaniya    siya        t{\USSmaller}um{\USGreater}awa        kahapon.\\
     \textsc{dat} \textsc{3sg.obl}    \textsc{3sg.nom}    <\textsc{av.rls}>laugh    yesterday\\
\glt `It was at him that she laughed yesterday.'
\end{xlist}
\end{exe}

\noindent Another inversion construction is the \isi{nominative} or \textit{ang-}inversion, in which the PSA is put in sentence‐initial position followed by a nominalized (\textit{ang}-marked) predicate phrase, yielding an equative structure like in (\ref{ang_inv2}). This construction is often translated by means of an \ili{English} \isi{cleft construction}, but can be shown to have a slightly different distribution \citep{latan16c} and a different syntactic structure \citep[e.g.][]{nagna07a}.

\begin{exe}
\ex\label{nom_inv1} \textbf{Nominative \textit{(ang)} inversion}
\begin{xlist}
\ex
\gll T{\USSmaller}um{\USGreater}a-tawa     siya.\\
     \textsc{<av.rls>ipfv}-laugh  \textsc{3sg.nom}\\
\glt `He was laughing.'
\ex\label{ang_inv2}
\gll Siya         ang     t{\USSmaller}um{\USGreater}a-tawa.\\
     \textsc{3sg.nom} \textsc{nom} \textsc{<av.rls>ipfv}-laugh\\
\glt `He was the one laughing.'
\end{xlist}
\end{exe}

\noindent Except for two instances of \textit{ang}-inversion (see (\ref{text4_part6}) in \sectref{sect:ct_bg_f} and (\ref{text5_part1}) in \sectref{sect:text_new}), speakers did not reproduce the last two inversion constructions in the study presented in this paper, although \cite{latan16c}
could show that for the Hunger Games corpus the \textit{ang}-inversion is more frequently used than the \ili{English} \textit{it}-cleft.

\section{The pragmatics of information structure and discourse structure}

\subsection{Information structure theory: Basic assumptions and terminology}\label{sect:is_theory}

In this and the following section, we will define the information-structural concepts and terminology used in the case studies of \sectref{sect:rats} from a \isi{discourse} perspective. At least two ways of describing \isi{information packaging} have been proposed in the literature; the topic-comment structure and the focus–background structure. \cite{Lambrecht1994} and others assume that the topic is the expression with respect to which the speaker aims to increase the hearer’s knowledge, the comment being the part which provides this knowledge. The focus-background packaging, on the other hand, is about the distinction of ``non-presupposed" vs.\ ``presupposed" material.\footnote{This usage of the term \textit{presupposed}, which can apply to denotations of non-clausal expressions and which is closely related to the notions \textit{given} or \textit{backgrounded}, stems from \cite{chono71} and \cite{jacra72}. It is, however, in conflict with the use of the term in formal-pragmatic approaches, e.g.\ \cite{karla74,staro74}, which essentially define presuppositions as (abstract) propositional entities assumed to be part of the common ground, i.e.\ not necessarily as given. Compare also discussions in \cite[150f.]{Lambrecht1994}, \cite{wagmi12} or \cite[237f.]{riear15}. For a contemporary picture of \isi{presupposition} and other types of projective content see e.g.\ \cite{tonju13}.} This definition is close in spirit, though not in terminology, to \cite{riear18a} or \cite{riear18b}, building on the tradition of Alternative Semantics \citep{rooma92} and theories of \textit{questions under discussion (\isi{QUD})} \citep{bueda03,bueda16,beada08,robcr12}. It is assumed that focus is the answer to the (current) \isi{QUD}, while the background is the content used when formulating the \isi{QUD}. Elements of the background that are not topical (here, interpreted as not referential, cf. \citealt{jacjo01}) are sometimes called \textit{tail} \citep[cf.][]{valen96}.

\begin{exe}
\ex \label{peter_topic1}
\begin{tabular}[t]{llllc}
\textit{Question:} & \multicolumn{4}{l}{What is Peter doing in the dark?}\\
\textit{Answer:} 
&\multicolumn{1}{c}{Peter} & is & dancing & in the dark. \\
&\textbf{[ Aboutness Topic ]} & \multicolumn{2}{l}{\textbf{[ Comment ]}} & \textbf{[\hfilll Tail \hfilll]}\\
& \multicolumn{2}{c}{\textbf{[\hfilll Background \hfilll]}}& \textbf{[ Focus ]} & \textbf{[  Background ]}\\
&\multicolumn{4}{l}{\textbf{[\hfilll Focus \hfill Domain \hfill \sq\hfilll]}}
\end{tabular}
\end{exe}

\noindent As for the information-structural markup, we choose the conventions demonstrated in (\ref{peter_topic2}), following  \cite{riear18b}: 

\begin{exe}
\ex\label{peter_topic2}
\begin{xlist}
\exi{Q\textsubscript{1}:} What is Peter doing in the dark?
\exi{> A\textsubscript{1}:} {[[Peter]\topic\ is [dancing]\focus\ in the dark]\sq}.
\end{xlist}
\end{exe}

\noindent (Q)uestions under discussion and their (A)nswers share the same index.  The indentation (>) signals the fact that answers are subordinate to their \isi{QUD} in \isi{discourse structure}, i.e.\ they stand in a parent-child relation, as symbolized in \figref{fig:simple1}.

\begin{figure}
\Tree [.{Q\textsubscript{1}}
        {A\textsubscript{1}}
      ]
\caption{Simple question-answer discourse}
\label{fig:simple1}
\end{figure}

\noindent The focus (F) is the part of the answer that corresponds to the question element of the \isi{QUD}. QUDs which are implicit are enclosed in curly brackets, while overt questions are represented without brackets. The \sq\ symbol \citep{rooma92} indicates a \isi{focus domain}, whose purpose it is to delimit the area comprising both the background and the focus. Focus domains ``match" \citep{bueda08} the \isi{QUD} regarding their background, and they also allow for the easy identification of mutually \isi{contrastive} assertions. All material inside the \isi{focus domain} that is not focal is part of the background. There is no separate label for the background itself but we identify referential expressions in the background, e.g.\ \textit{Peter} in example (\ref{peter_topic1}), as aboutness topics (T).\footnote{As for functional elements (e.g.\ the copula \textit{is}, determiners, or prepositions), we leave it open when they should become part of the focus and when they should not.}

A further category, which is not yet satisfactorily captured by the above grouping, is the \isi{contrastive topic} \textsc{(ct)}. Following \cite{bueda03}, a \isi{contrastive topic} is a hybrid category, which shares properties with both focus and background (or \isi{aboutness topic}). We assume that \isi{contrastive} topics can only occur in combination with a focus. They reflect a complex \isi{discourse} strategy of the kind depicted in (\ref{ct_f2}) or, more abstractly, in \figref{fig:ct_simple1}.

\newpage 
\begin{exe}
\ex\label{ct_f2}
\begin{xlist}
\exi{Q\textsubscript{1}:} \{Who is doing what in the dark?\}
\exi{> Q\textsubscript{{1.1}}:}  \{What is Peter doing in the dark?\}
\exi{{>}> A\textsubscript{{1.1}}:} {[[Peter]\ctopic\ is [dancing]\focus\ in the dark]\sq}.
\exi{> Q\textsubscript{{1.2}}:} \{What is Anna doing in the dark?\}
\exi{{>}> A\textsubscript{{1.2}}:} {[[Anna]\ctopic\ is [eating]\focus\ in the dark]\sq}.
\end{xlist}
\end{exe}

\begin{figure}
\begin{forest} %\Tree breaks here
[{Q\textsubscript{1}}
    [{Q\textsubscript{{1.1}}}
      [{A\textsubscript{{1.1}}}]
    ]
    [{Q\textsubscript{{1.2}}}
      [{A\textsubscript{{1.2}}}]
    ]
  ] 
\end{forest} 
\caption{Question-subquestion discourse}
\label{fig:ct_simple1}
\end{figure}

\noindent The complex strategy consists in the fact that a question with two question elements (here: \textit{who, what}) are answered by a sequence of answers to subquestions (indicated Q\textsubscript{{1.1}}, Q\textsubscript{{1.2}}). Now, the expressions \textit{Peter/Anna} -- which are \isi{backgrounded} with respect to the subquestions but focal with respect to the main question Q\textsubscript{1} -- are \isi{contrastive} topics. Note that \isi{contrastive} topics, as in (\ref{ct_f1}), can function as aboutness topics (i.e.\ indicate a \isi{discourse} \isi{referent}) but they need not. Throughout the languages of world, it is the \isi{contrastive} function of expressions that leads to stronger deviations from the default sentence realization than the aboutness function. In other words, it seems to be more important to signal that a change is expected or currently happening than to signal that there is a mere topical continuity. For that reason we expect \isi{contrastive} topics to trigger more dramatic changes with regard to \isi{constituent order} and \isi{prosody} than (non-\isi{contrastive}) aboutness topics.

\subsection{Discourse structure trees, QUDs and well‐formedness conditions}\label{sect:principles}

Recently, \cite{riear18b} (see also \citealt{reyuw16}, \citealt{riear18a}) have proposed a procedure for the identification of implicit questions under discussion (QUDs) in textual data. We will apply their methodology in the data analyses of \ili{Tagalog} presented in \sectref{sect:rats}. The method implements an insight, going back to at least \cite{stuch89,kupja95,ginjo96a} and \cite{robcr12}, that the assertions made in a text can actually be thought of as answers to implicit questions, much like question-answer sequences in spoken dialogues. While in much of the previous work, QUDs have remained an object of theoretical investigation, \cite{riear18b} develop practical linguistic annotation guidelines, which are demonstrated on the basis of \ili{French}, \ili{German} and \ili{English} corpus data.\footnote{See also \citetv{riear18c} for an application of the QUD-tree method to \ili{Sumbawa} (\ili{Austronesian}, \ili{Malayo-Polynesian}).} Since the information structure of an \isi{assertion} is definable relative to its \isi{QUD}, the benefit of enhancing a text with pragmatically defined implicit QUDs is that we gain access to the information structure of each clausal unit. This provides us with potential access to a large amount of data, which can be used to study the morphosyntactic marking of information structure in any language of interest. The method is more flexible, faster and cheaper than collecting data by means of experimental techniques, and comes with the additional advantage that the data under investigation are potentially more natural than, for instance, artificially produced question-answer responses.
Without going into details, the reconstruction of the implicit QUDs of a text is enabled because it is constrained by a number of pragmatic principles derived and adapted from the literature on information structure; in particular \cite{rooma92,schro99,bueda03} and \cite{bueda08}:
 
\begin{itemize}
\item For any \isi{assertion} $A$ identified in a text, its immediately dominating question $Q$ must be directed at one of its constituents (\textsc{Q-A-Con\-gru\-ence}); i.e.\ QUDs which do not target any expression in the \isi{assertion} are forbidden.

\item The principle of \textsc{Q-Givenness} says that $Q$ can only consist of material that is salient at the point when $A$ occurs, where \textit{salient} means that the content is already active in the mind of the reader before $A$ is processed. Typically, activation results from previous mention. In other words, material that is given in the \isi{discourse} counts as salient and can be used to formulate $Q$.

\item The \textsc{Q-Givenness} constraint is complemented by the principle \textsc{Maximize-Q-Ana\-pho\-ri\-ci\-ty}, which says that all the content that is given in \isi{assertion} $A$ should be reflected in the \isi{QUD}; i.e.\ $A$ should have a maximal background and a minimal focus.
\end{itemize}

\noindent For instance, in (\ref{rats_info1}) the appropriate \isi{QUD} for A\textsubscript{1} must be the one indicated by Q\textsubscript{1}, because it is this question which contains only given material \textit{(rats)} and maximizes the background of A\textsubscript{1}. By contrast, the reader is encouraged to check that other questions (e.g.\ \textit{Is stress unhealthy? What about cats? What happened?}) violate one or several of the constraints mentioned above. 

\begin{exe}
\ex\label{rats_info1}
\begin{xlist}
\exi{A\textsubscript{0}:} The life of rats is stressful and dangerous.
\exi{Q\textsubscript{1}:} \{What about rats?\} 
\exi{> A\textsubscript{1}:} {[[They]\topic\ [get chased by cats]\focus]\sq}.
\end{xlist}
\end{exe}

\noindent The final constraint, \textsc{Parallelism}, is designed to handle \isi{contrastive} information in a text. When two assertions $A'$ and $A''$ share semantic content while their remaining parts are interpreted as alternatives, then there is a common \isi{QUD}, in which the shared material is reflected. In specific cases, \textsc{Parallelism} will override \textsc{Q-Givenness}, which means that the \isi{QUD} of two (or more) parallel assertions may contain some discourse-new material if it is (semantically) contained in both assertions. As an example, consider the sequence in (\ref{rats_parallel1}). The information-structurally relevant \isi{QUD} is Q\textsubscript{{1.1}} (a subquestion of Q\textsubscript{1}) which includes the discourse-new (but semantically shared) material \textit{chase}.\footnote{This example demonstrates that it makes sense to distinguish between \textit{topical} reference phrases, e.g.\ \textit{[rats]\topic}, and larger, non-referential \textit{backgrounds} that may contain more material than just the \isi{aboutness topic}, e.g.\ the phrase \textit{[chase [rats]\topic]}.}

\begin{exe}
\ex\label{rats_parallel1}
\begin{xlist}
\exi{A\textsubscript{0}:} The life of rats is stressful and dangerous.
\exi{Q\textsubscript{1}:} \{What about rats?\}
\exi{> Q\textsubscript{{1.1}}:} \{Who chases rats?\}
\exi{{>}> A\textsubscript{{1.1'}}:} {[[Cats]\focus\ chase [rats]\topic]\sq}
\exi{{>}> A\textsubscript{{1.1''}}:} {and [[dogs]\focus\ go after [them]\topic]\sq\ as well.}
\end{xlist}
\end{exe}

\noindent Finally, example (\ref{rats_nai1}) introduces so-called non-at-issue (\textsc{nai}) material \citep{potch05}, which we define as linguistically optional expressions or, more precisely, as discourse-new material which is syntactically and semantically independent and which does not contribute in a direct manner to answering the current \isi{QUD}. Typical \textsc{nai} expressions are, for instance, appositions, non-restrictive relative clauses, speaker-oriented adverbs, evidentials or adjunct clauses like in (\ref{rats_nai1}).

\begin{exe}
\ex\label{rats_nai1}
\begin{xlist}
\exi{Q\textsubscript{0}:} \{What do cats do?\}
\exi{> A\textsubscript{0}:} {[When they are in the mood,]\nai\ [[cats]\topic\ [chase big rats]\focus]\sq.}
\end{xlist}
\end{exe}

\noindent We will simply ignore non-at-issue material when it occurs at the beginning or in the middle of a \isi{clause}.  Non-at-issue material at the end of a \isi{clause} will be treated as a new, independent \isi{assertion}. Finally, note that \isi{discourse} markers (e.g.\ \textit{and, or, but, although}) or \isi{discourse} particles (\textit{also, as well, only, even}) are not marked as \textsc{nai} but simply left unannotated; compare A\textsubscript{{1.1''}} in (\ref{rats_parallel1}).

\section{Case study `The unhappy rats’}\label{sect:rats}

\subsection{QUD approach to `The unhappy rats’}\label{sect:rats_texte}

Our case study is based on data elicited in Manila. Four speakers of \ili{Tagalog} (three of which only speak \ili{Tagalog} and no other Philippine language, and one who speaks \ili{Palawan} and \ili{Tagalog}) were asked to freely translate five short texts that differ in the \isi{givenness} of the \isi{actor argument}, the \isi{undergoer} argument, the verb, and in the questions under discussion answered by the target sentence. The target sentences themselves vary slightly in their formulation in order to make the discourses sound more natural. In the following overview of the material we focus on the \isi{givenness} of the arguments in the target sentences. The \isi{discourse structure} of the texts will be worked out in detail in the next section.

In Text 1, the actor \textit{cats} is the \isi{discourse topic}. The verb and the \isi{undergoer} provide the new information.

\begin{exe}
\exi{1.}\textit{Cats are silly creatures with nothing but nonsense on their minds. They climb up on curtains, bring home mice. \grey{Cats also chase and catch big rats,} when they get a chance. Who wants to have a big rat in their house?}
\end{exe}

\noindent In Text 2, the \isi{undergoer} \textit{rats} is the \isi{discourse topic}, and there is a \isi{narrow contrastive focus} on the actor \textit{cats}.

\begin{exe}
\exi{2.}\textit{It is not only wolves and foxes that threaten rats and catch them. \grey{Cats also catch rats} and eat them afterwards.}
\end{exe}

\noindent In Text 3, the \isi{undergoer} \textit{rats} is the \isi{discourse topic}. The actor \textit{cats} and their activity with respect to the \isi{undergoer} is the new information.

\begin{exe}
\exi{3.} \textit{Rats live stressful and dangerous lives. The noise of the traffic makes them nervous and sick. Dogs chase them.  And also \grey{(our domestic) cats catch and kill rats,} when they get the chance.}
\end{exe}

\noindent In Text 4, the \isi{discourse topic} is the cruel laws of nature.  All sentences are about predators catching certain animals. Thus, the verb combination \textit{chase and catch} is given, while the actor \textit{cats} and \isi{undergoer} \textit{rats} are new/\isi{contrastive}.

\begin{exe}
\exi{4.} \textit{Life in the wilderness is pretty cruel. Lions catch antelopes, sharks catch tunafish and happen to get caught and killed by humans themselves. Even here in the city  these cruel laws of nature can be observed. \grey{(Our domestic) cats also chase and catch rats,} and  some also bring them home to continue playing with the bleeding creature.} 
\end{exe}

\noindent In Text 5, the  \isi{discourse topic} is the unhappiness and violence that the narrator witnesses. All follow-up sentences are all-new sentences that elaborate on what kind of violence the narrator observes.

\begin{exe}
\exi{5.} \textit{When I look out of the window, I see only unhappiness and violence.  Dogs chase hens and make them lose their feathers. Old bitter women scream at children and make them cry.  And also \grey{(our domestic) cats catch and kill innocent rats,} when no one is looking.} 
\end{exe}

\noindent A PhD student and one of the authors presented the \ili{Tagalog} translations to nine native speakers of \ili{Tagalog} in Manila in a private setting. The consultants had different educational backgrounds and were aged 20--63. They were not aware of the original \ili{English} texts and were asked to read, evaluate, correct and improve the \ili{Tagalog} versions of the five texts. They were also asked which translations they liked better. Interestingly, the participants never corrected sentence structure or \isi{voice selection}, but only vocabulary choices. When asked explicitly about variants that could be found in different texts with respect to syntax and \isi{voice}, they stated their preferences, but none of the translations was rejected as awkward. We therefore consider the translations we got as acceptable and natural \ili{Tagalog}.

The nature of the texts led to translations with a lot of parallelisms and the occurrence of the particle \textit{rin} `also' in the target sentence \textit{Cats also catch rats,} regardless of whether or not this particle had been given in the \ili{English} original. As pointed out by \cite{krima99}, particles like \textit{also} tend to have an associated constituent, which is often the \isi{contrastive topic}.

The overall goal of the study was to find out how far the difference in textual \isi{givenness} of the arguments would influence \isi{voice selection} and construction choice in the target sentence.  Recall that under the \isi{QUD} approach \citep{riear18b} the crucial information-structural classes distinguished are focus, background (including \isi{aboutness topic}), \isi{contrastive topic}, and non-at-issue material, while a differentiation between new and \isi{contrastive focus} is not made.
 
In the following sections we will provide an analysis of the \isi{discourse} structures and the implicit questions under discussion giving rise to comprehensive information-struc\-tu\-ral analyses.  Remember that the discourses shown in \sectref{sect:text1}--\sectref{sect:text_new} differ from the ones introduced above in this section, since we present the free \ili{Tagalog} translations of the original texts as well as their re-translations into \ili{English}. Based on the analyses we can check whether our expectations regarding the coding of the target sentence and the actual coding choices by native speakers match.

\subsection{Text 1 -- Intended: Actor (topical) Verb (focal) Undergoer (focal)}\label{sect:text1}

\begin{figure}
\framebox{
\small
\begin{tikzpicture}
\tikzset{inner xsep=0em,level distance=5em,sibling distance=0em,every tree node/.style={text width=8em, text centered}}
\Tree [.{Q\textsubscript{0}: \{What is the way things are?\}}
        [.{Q\textsubscript{{0.1}}: \{What do different animals do?\} }
          [.{Q\textsubscript{{0.1.1}}: \{What do cats do?\}}
 		    {A\textsubscript{{0.1.1}}: Cats do a lot of things without sense.}
              [.{Q\textsubscript{{0.1.1.1}}: \{What exactly do they do?\}}
                {A\textsubscript{{0.1.1.1'}}: They climb up on curtains,}
                {A\textsubscript{{0.1.1.1''}}: and bring home mice.}
                {\textbf{A\textsubscript{{0.1.1.1'''}}: Cats also catch big mice (= rats)} (when there is an opportunity).}
              ]
          ]
        ]
      ]
\end{tikzpicture}
}
\caption{Discourse structure of Text~1 (free Tagalog formulation, translated back into English)}
\label{fig:text1_analysis}
\end{figure}

\noindent The first short text deals with the life of cats and the silly things they do. The QUD-tree analysis is shown in \figref{fig:text1_analysis} (target sentence in boldface).
 The actor \textit{cat} was thus given, and intended was focus on the verb and the \isi{undergoer} \textit{rats}, i.e.\ \textit{[[cats]\topic\ also [catch rats]\focus]\sq}. The text was
conceived in a way that no information-structural \isi{prominence} in the sense of (\ref{is_prom}) for either argument needed to be marked morphosyntactically. Based on Hypothesis (iii) developed in \sectref{sect:pat_prom}, we predicted no special syntactic structure for the target sentence in this case, i.e.\ only the normal predicate-initial structure. But we expected \isi{actor voice} to be the preferred choice for this scenario because the undergoers are non-specific and the verbs are not undergoer-oriented. This is indeed what we found, as A\textsubscript{{0.1.1.1'''}}, discussed below in (\ref{text1_part3}), shows. While we only give the translation of one consultant here, all other consultants chose the same form, i.e.\ a predicate-initial sentence with \isi{actor voice}.  

According to \cite{robcr12}, any \isi{discourse} addresses the so-called ``Big Question" Q\textsubscript{0}: \textit{\{What is the way things are?\}}  Many actual discourses, however, start \textit{in medias res} \citep{firja92}, i.e.\ they use grammatical means (like \isi{presupposition} triggers or non-default \isi{constituent order}) to express that certain information should be accommodated (i.e.\ treated as if it were given), which boils down to answering a more specific subquestion (indicated in our example text by Q\textsubscript{{0.1}} and Q\textsubscript{{0.1.1}}\footnote{Note that the sub-numbering convention indicates that there is an entailment relation between the questions. For instance, every answer to Q\textsubscript{{0.1.1}} is, at the same time, an answer to Q\textsubscript{{0.1}} and, in turn, to Q\textsubscript{0}.}). Since such grammatical means differ from language to language, we cannot be sure, at the outset, that our QUD-analysis for the initial sentence is correct. It is generally advisable not to draw any crucial conclusions from discourse-initial sentences (which we don’t). In this example, we merely propose, without proving it, that the speaker has chosen to start a \isi{contrastive} \isi{discourse} strategy, analogous to the one in (\ref{ct_f2}), \sectref{sect:is_theory}, with the goal to provide information about different animals. We suggest therefore to analyse the expression \textit{cats} in A\textsubscript{{0.1.1}} of (\ref{text1_part1}) as a \isi{contrastive topic}. Combinations of \textsc{ct} and \textsc{f} can be coded via \textit{ay}-inversion in \ili{Tagalog}, as shown in example (\ref{ct_f1}) and suggested in Hypothesis (i), cf.\ \cite{latan17}.

\begin{exe}
\ex\label{text1_part1}
\begin{xlist}
\exi{Q\textsubscript{0}:} \{What is the way things are?\}
\exi{> Q\textsubscript{{0.1}}:} \{What do different animals do?\} 
\exi{{>}> Q\textsubscript{{0.1.1}}:} \{What do cats do?\}
\end{xlist}
\begin{xlist}[{>}{>}> A\textsubscript{{0.1.1}}:]
\exi{{>}{>}> A\textsubscript{{0.1.1}}:}
\gll {{\ob}{\ob}Ang} mga {pusa{\cb}\ctopic} ay {{\ob}g{\USSmaller}um{\USGreater}a-gawa} ng mga wala-ng {kabuluha-ng} {bagay{\cb}\focus{\cb}\sq}.\\
\textsc{\void{[[}nom} \textsc{pl} cat \textsc{inv} \textsc{\void{[}<av.rls>ipfv}-do \textsc{gen} \textsc{pl} \textsc{neg-lk}   sense-\textsc{lk}  thing\\
\glt `Cats do a lot of things without sense.'
\end{xlist}
\end{exe}

\noindent Assertion A\textsubscript{{0.1.1}} is then elaborated by use of three partial answers to the question in Q\textsubscript{{0.1.1.1}} in (\ref{text1_part2}). The three parallel answers, hence, all carry focus on the predicate and \isi{undergoer}. Given that the actor is topical and the \isi{undergoer} is focal, i.e.\ given that the arguments have their IS-default values, Hypothesis (iii) delegates the \isi{voice selection} to the level of \isi{referentiality}. As the \isi{referent} of the \isi{undergoer} argument is non-specific, like in example (\ref{write_av1}) in the introduction, \isi{actor voice} is expected to be preferred, and this is also what we find in the translations of all four consultants. Note that the sentences in (\ref{text1_part2}) and (\ref{text1_part3}) have default \isi{constituent order}. The reason why the PSA is not in final position here is its occurrence as a (\isi{clitic}) \isi{pronoun}.  

\begin{exe}
\ex\label{text1_part2}
\begin{xlist}
\exi{{>}{>}> Q\textsubscript{{0.1.1.1}}:\hfill} \{What exactly do they do?\}
\end{xlist}
\begin{xlist}[{>}>{>}> A\textsubscript{{0.1.1.1'''}}:]
\exi{{>}>{>}> A\textsubscript{{0.1.1.1'}}:\hfill}
\gll {{\ob}{\ob}Um-a-akyat{\cb}\focus} {{\ob}sila{\cb}\topic} {{\ob}sa} mga {kurtina{\cb}\focus{\cb}\sq}\\
\textsc{\void{[[}av-ipfv}-climb \textsc{\void{[}3pl.nom} \textsc{\void{[}dat} \textsc{pl} curtain\\
\glt `They climb up on curtains,'
\exi{{>}>{>}> A\textsubscript{{0.1.1.1''}}:\hfill}
\gll at    {{\ob}{\ob}nag-da-dala}    ng mga daga    sa {bahay{\cb}\focus{\cb}\sq}.\\
 and \textsc{\void{[[}av.rls-ipfv}-carry \textsc{gen} \textsc{pl} mouse   \textsc{dat} house\\
\glt `and bring home mice.'
\end{xlist}
\end{exe}
\begin{exe}
\ex\label{text1_part3}\textbf{Target sentence:}
\begin{xlist}[{>}>{>}> A\textsubscript{{0.1.1.1''':}}]
\exi{{>}>{>}> A\textsubscript{{0.1.1.1''':}}}
\gll {{\ob}Kung} may {pagkakataon,{\cb}\nai} {{\ob}{\ob}nang-hu-huli{\cb}\focus} rin {{\ob}sila{\cb}\topic} {{\ob}ng} {malaki-ng} {daga{\cb}\focus{\cb}\sq}.\\
{\void{[}if} \textsc{exist} opportunity, \textsc{\void{[[}av.rls-ipfv}-catch also \textsc{\void{[}3pl.nom} \textsc{\void{[}gen} big-\textsc{lk} mouse\\
\glt `Cats also catch big mice (= rats), when there is an opportunity.' 
\end{xlist}
\end{exe}

\noindent Note that, in the final \isi{assertion} A\textsubscript{{0.1.1.1'''}} in (\ref{text1_part3}) the preceding \textit{if}-sentence, labelled as not at issue, does not answer the question under discussion Q\textsubscript{{0.1.1.1}}, but rather expresses the relevance condition for the truth of the sentence containing the actual at-issue content. The target sentence \textit{cats also catch rats} has a parallel structure to the two preceding sentences that answer the same question under discussion. All sentences have the same basic structure, i.e.\ they are predicate-initial and show \isi{actor voice}. Note that the focus-sensitive particle \textit{rin} appears directly in the postverbal position in A\textsubscript{{0.1.1.1'''}}. 

\subsection{Text 2 -- Intended: Narrow Actor Focus, i.e.\ Actor (focal) Verb (backgrounded) Undergoer (topical); Partly realized: Actor (contrastive topic) Verb (focus) Undergoer (topical)}\label{sect:narrow_actor_focus}

\begin{figure}
\framebox{
\begin{tikzpicture}
\tikzset{inner xsep=0em,level distance=3em,sibling distance=0em,every tree node/.style={text width=15em, text centered}}
\Tree [.{Q\textsubscript{0}: \{What is the way things are?\}}
        [.{Q\textsubscript{{0.1}}: \{Who catches rats?\}}
          {A\textsubscript{{0.1'}}: It is not only bears and wolves that threaten and catch rats.}
          {\textbf{A\textsubscript{{0.1''}}: Cats catch them.}}
        ]
      ]
\end{tikzpicture}
}
\caption{Discourse structure of free formulation based on Text~2}
\label{fig:text2_analysis}
\end{figure}

\noindent The second short text\footnote{For unknown reasons, the consultant changed the names of the animals from the original text.} (\isi{discourse structure} given in \figref{fig:text2_analysis}) was intended to yield a \isi{narrow focus} on the \isi{actor argument} in the target sentence \textit{[[Cats]\focus\ also catch [rats]\topic]\sq} and, thus, represent one of the marked constructions mentioned in Hypotheses (iv) and (v), i.e.\ either Verb\av\ \isi{Undergoer}\topic\  Actor\focus\ \textit{or} Actor\focus\ +ay Verb\uv\ \isi{Undergoer}\topic, because both arguments are information-structurally prominent according to (\ref{is_prom}). In order to force a context that would yield a \isi{narrow actor focus} question as the natural implicit question under discussion for the target sentence, we provided a preceding sentence containing the focus sensitive particle \textit{l(am)ang} `only' with scope over those actor arguments that the actor \textit{cat} in the target sentence was supposed to be contrasted with. The negation of the exhaustive particle in the first sentence was meant to be an indicator of the ensuing parallelism. In order to express the (negated) exhaustive \isi{narrow focus} on the \isi{actor argument} expressed in the context sentence, the \textit{ang}-inversion construction (compare (\ref{ang_inv2}) in \sectref{sect:is_tagalog}) was chosen by all \ili{Tagalog} translators, as shown in (\ref{text2_part1}).

\begin{exe}
\ex\label{text2_part1}
\begin{xlist}
\exi{Q\textsubscript{0}:} \{What is the way things are?\}
\exi{> Q\textsubscript{{0.1}}:} \{Who catches rats?\}
\end{xlist}
\begin{xlist}[{>}> A\textsubscript{{0.1'}}:]
\exi{{>}> A\textsubscript{{0.1'}}:}
\gll {{\ob}Hindi} lang {{\ob}ang} mga oso at {lobo{\cb}\focus} ang na-nakot    at nag-hu-huli {{\ob}ng} mga {daga{\cb}\topic{\cb}\sq}.\\
\textsc{\void{[}neg} only \textsc{\void{[}nom} \textsc{pl} bear and wolf \textsc{nom} \textsc{av.rls}.scare and \textsc{av.rls-ipfv}-catch \textsc{\void{[}gen} \textsc{pl} rat\\    
\glt `It is not only bears and wolves that threaten and catch rats.'
\end{xlist}
\end{exe}
 
\noindent Subsequently, speakers did not continue with another \textit{ang}-inversion but chose \textit{ay}-inver\-sion for the target sentence. As for \textit{ay}-inverted \isi{narrow actor focus}, Hypothesis (v) made us expect \isi{undergoer voice} to mark the 
\isi{topicality} of the \isi{undergoer}. Indeed, two consultants offered this construction, which is given in (\ref{text2_part3}).\footnote{A predicate-initial construction as in (\ref{martinez_b}), with \isi{undergoer voice} and the genitive-marked actor in sentence-final, i.e.\ a prosodically prominent position, was rejected as ``not good" by all consultants, i.e.\ it was neither offered in the translations nor accepted as a possible option, when we asked the speakers in Manila about this. This adds to our suspicion that the construction in question does not express \isi{narrow contrastive focus}.} The two other consultants suggested the construction in (\ref{text2_part2}), i.e.\ \textit{ay}-inversion of the actor combined with \isi{actor voice}. 

\begin{exe}
\ex\label{text2_part3}\textbf{Target sentence:}
\begin{xlist}[{>}> A\textsubscript{{0.1''}}:]
\exi{{>}> A\textsubscript{{0.1''}}:} 
\gll {{\ob}{\ob}Ang} mga {pusa{\cb}\focus} din ay h{\USSmaller}in{\USGreater}u-huli {{\ob}sila{\cb}\topic{\cb}\sq}.\\
\textsc{\void{[[}nom} \textsc{pl} cat also \textsc{inv} \textsc{<uv.rls>-ipfv}-catch \textsc{\void{[}3pl.nom}\\ 
\glt `Cats also catch them.'
\end{xlist}
\end{exe}

\begin{exe}
\ex\label{text2_part2}\textbf{Alternative realization of target sentence:}
\begin{xlist}
\exi{{>}> Q\textsubscript{1}:}{\{Who does what to rats?\}}
\exi{{>}{>}> Q\textsubscript{{1.1}}}{\{What do cats do to rats?\}}
\end{xlist}
\begin{xlist}[{>}>{>}> A\textsubscript{{1.1}}:]
\exi{{>}>{>}> A\textsubscript{{1.1}}:}
\gll {{\ob}{\ob}Ang} mga {pusa{\cb}\ctopic} {\USOParen}rin{\USCParen} ay {{\ob}nang-hu-huli} at {nag-ka-kain{\cb}\focus} {{\ob}ng} mga {daga{\cb}\topic{\cb}\sq}.\\
\textsc{\void{[[}nom} \textsc{pl} cat (also) \textsc{inv} \textsc{\void{[}av.rls-ipfv}-catch and \textsc{av.rls-ipfv}-eat \textsc{\void{[}gen} \textsc{pl} rat\\
\glt `Cats (also) catch and eat rats.'
\end{xlist}
\end{exe}

\noindent Note that the target sentence A\textsubscript{{1.1}} in (\ref{text2_part2}), which has a different label than the one in (\ref{text2_part3}) due to a very different \isi{discourse structure},  contains a given verb \textit{catch} and a new verb \textit{eat}. It is quite likely that the way Text~2 was formulated had the (unintended) effect that the consultants construed the verb coordination as being contrasted against the previous combination \textit{threaten and catch}. Hence, we assume, as indicated in (\ref{text2_part2}), that a \textsc{ct+f} structure was chosen with focus on the verb complex. This also leads us to diversify our inventory of hypotheses once more and to add a slight change to Hypothesis (v), restricting it to \isi{backgrounded} verbs.

\hypo{v}{Actor\focus\ +ay Verb\bg\ \isi{Undergoer}\topic\ $\rightarrow$ \isi{Undergoer} Voice\hfill\textit{(final)}} 

\hypo{vi}{Actor\ctopic\ +ay Verb\focus\ \isi{Undergoer}\topic\ $\rightarrow$ Actor Voice} 

\noindent While it may seem ad hoc to formulate a new hypothesis merely on the basis of a somewhat unclear example like (\ref{text2_part2}), we will come across a very similar example in the next section, which seems to confirm that Hypothesis (vi) is on the right track.

Most of the \ili{Tagalog} target sentences for Text~2 did not contain a word for \textit{also}, which was contained in the original \ili{English} text; somehow consultants seemed to feel that the \textit{ay}-construction already conveyed an additive focus reading. Only one consultant paid heed to the focus sensitive particle in her translation and placed \textit{rin} `also' right after the actor in one instantiation of (\ref{text2_part2}). As we have seen above and see again here, the positioning of \textit{rin} indicates its scope. If the verb is part of the scope, \textit{rin} appears after verb. If only the actor is in its scope, it appears right after the actor. 

  
\subsection{Text 3 -- Intended: Actor (focal) Verb (focal) Undergoer (topical); Result: Actor (contrastive topic) Verb (focal) Undergoer (topical)}

\begin{figure}
\framebox{
\footnotesize
\begin{tikzpicture}
\tikzset{inner xsep=0em,level distance=5em,sibling distance=0em,every tree node/.style={text width=6em,text centered}}
\Tree [.{Q\textsubscript{0}: \{What is the world like?\}}
        [.{Q\textsubscript{{0.1}}: \{What about the life of rats?\}}
          {A\textsubscript{{0.1}}: The life of rats is stressful and dangerous.}
          [.{Q\textsubscript{1}: \{What about rats?\}}
            {A\textsubscript{{1'}}: They are made nervous}
            {A\textsubscript{{1''}}: and become sick}
            [.{Q\textsubscript{2}: \{By what are they made sick?\}}
              {A\textsubscript{2}: due to the noise of the traffic.}
            ]
            [.{Q\textsubscript{3}: \{Who does what to rats?\}}
              [.{Q\textsubscript{{3.1}}: \{What do cats do to rats?\}}
                {\textbf{A\textsubscript{{3.1}}: And also our domestic cats catch rats}}
                [.{Q\textsubscript{4}: \{When do they do this?\}}
                  {A\textsubscript{4}: when they get the chance.}
                ]
              ]
            ]
          ]
        ]
      ]
\end{tikzpicture}
}
\caption{Discourse structure based on Text~3}
\label{fig:text3_analysis}
\end{figure}

\noindent Next, we wanted to find out whether a truly new verb would make a difference for the construction chosen, so we construed a text in which both the actor and the verb were discourse-new. This text (\isi{discourse structure} shown in \figref{fig:text3_analysis}) is about the life of rats, i.e.\ it is about the \isi{undergoer} of the target sentence.  In the sentences preceding the target sentence there is a clear focus on the events that affect rats. The translation provided by one of our consultants contains simple predicate-initial structures. This is in line with our expectation regarding \isi{predicate focus} sentences.

\begin{exe}
\ex \label{text3_part1}
\begin{xlist}
\exi{Q\textsubscript{0}:} \{What is the world like?\}
\exi{> Q\textsubscript{{0.1}}:} \{What about the life of rats?\}
\end{xlist}
\begin{xlist}[{>}> A\textsubscript{{0.1}}:]
\exi{{>}> A\textsubscript{{0.1}}:}
\gll {{\ob}{\ob}Delikado} at {ma-stress{\cb}\focus} {{\ob}ang} buhay  ng mga {daga{\cb}\topic{\cb}\sq}.\\
{\void{[}Dangerous} and stressful \textsc{\void{[}nom} life \textsc{gen} \textsc{pl} rats\\
\glt `The life of rats is stressful and dangerous.'
\end{xlist}
\end{exe}

\begin{exe}
\ex \label{text3_part2}
\begin{xlist}[{>}> A\textsubscript{{0.1}}:]
\exi{{>}> Q\textsubscript{1}:\hfill} \{What about rats?\}
\end{xlist}
\begin{xlist}[{>}{>}> A\textsubscript{{1''}}:]
\exi{{>}{>}> A\textsubscript{{1'}}:\hfill}
\gll {{\ob}{\ob}Ni-ni-nerbyos{\cb}\focus} {{\ob}sila{\cb}\topic{\cb}\sq}\\
\textsc{\void{[[}uv.rls-ipfv}-nervous \textsc{\void{[}3pl.nom}\\
\glt `They are made nervous'
\exi{{>}{>}> A\textsubscript{{1''}}:}
\gll at {{\ob}{\ob}nag-kakasakit{\cb}\focus{\cb}\sq}\\   
and \textsc{\void{[[}av.rls}-become.sick\\
\glt `and become sick,'
\end{xlist}
\end{exe}

\begin{exe}
\ex \label{text3_part3}
\begin{xlist}[{>}{>}> A\textsubscript{{1''}}:]
\exi{{>}{>}> Q\textsubscript{2}:\hfill} \{By what are they made sick?\}
\end{xlist}
\begin{xlist}[{>}>{>}> A\textsubscript{2}:]
\exi{{>}>{>}> A\textsubscript{2}:}
\gll {{\ob}{\ob}gawa} ng ingay ng {trapik{\cb}\focus{\cb}\sq}.\\
{\void{[[}done} \textsc{gen} noise \textsc{gen} traffic\\
\glt `due to the noise of the traffic.'\\
Orig: `The noise of the traffic makes them nervous and sick.’
\end{xlist}
\end{exe}

\noindent Interestingly, the consultant combines an \isi{undergoer voice} verb and an \isi{actor voice} verb in (\ref{text3_part2}), realizing the affector-actor within an optional adjunct phrase, which we treat as a separate information unit at the end of the sentence. We therefore have narrow (\isi{contrastive}) \isi{predicate focus} on the verbs in (\ref{text3_part2}) and a secondary focus on the actor phrase, which is not at issue with respect to Q\textsubscript{1} but at issue with respect to the separate subquestion Q\textsubscript{2} in (\ref{text3_part3}). Note that the way the text was construed, we had expected to obtain a parallel structure between two actors (or affectors), the traffic and the cats. However, since the speaker chose to demote the traffic to a peripheral adjunct and used the intransitive verb \textit{to fall sick}, a non-parallel \isi{discourse structure} resulted. 

Similarly to the narrow-actor-focus context, i.e.\ example (\ref{text2_part2}) in \sectref{sect:narrow_actor_focus}, the actor phrase \textit{our domestic cats} in the target sentence in (\ref{text3_part4}) was not realized at the end of the sentence (hence, Hypothesis (iv) is not applicable), but in sentence-initial position before the inversion marker \textit{ay} and, once again, \isi{actor voice} instead of \isi{undergoer voice} was chosen, despite the given \isi{undergoer}.

\begin{exe}
\ex\label{text3_part4}\textbf{Target sentence:}
\begin{xlist}
\exi{{>}{>}> Q\textsubscript{3}:} \{Who does what to rats?\}
\exi{{>}>{>}> Q\textsubscript{{3.1}}:} \{What do cats do to rats?\}
\end{xlist}
\begin{xlist}[{>}>{>}{>}> A\textsubscript{{3.1}}:]
\exi{{>}>{>}{>}> A\textsubscript{{3.1}}:}
\gll At {{\ob}{\ob}ang} ating mga alagang {pusa{\cb}\ctopic} ay {{\ob}h{\USSmaller}um{\USGreater}u-huli{\cb}\focus} rin {{\ob}ng} {daga{\cb}\topic{\cb}\sq}\\
and \textsc{\void{[[}nom} our \textsc{pl} pet cat \textsc{inv} \textsc{\void{[}<av.rls>ipfv}-catch also \textsc{\void{[}gen} rat\\
\glt `And also our domestic cats catch rats,'
\end{xlist}
\end{exe}
\begin{exe}
\ex\label{text3_part5}
\begin{xlist}[{>}>{>}{>}> A\textsubscript{{3.1}}:]
\exi{{>}>{>}{>}> Q\textsubscript{4}:\hfill} \{When do they do this?\}
\end{xlist}
\begin{xlist}[${>}>{>}>{>}>$ A\textsubscript{4}:]
\exi{${>}>{>}>{>}>$ A\textsubscript{4}:}
\gll {{\ob}{\ob}kapag} may {pagkakataon{\cb}\focus{\cb}\sq}.\\
{\void{[[}if} exist chance\\
\glt `when they get the chance.'
\end{xlist}
\end{exe}

\noindent We cannot claim that the speakers simply gave the exact same structure to this utterance as to the \isi{narrow actor focus} construction in (\ref{text2_part2}), because the particle \textit{rin} appears after the verb, not after the actor. Otherwise, however, the syntactic structure and \isi{voice} are identical and, hence, in accordance with Hypothesis (vi). As witnessed above, the \textit{ay}-inversion seems to express the first half of a \isi{contrastive} \isi{discourse} strategy, which is mirrored in the assumption of Q\textsubscript{3} and Q\textsubscript{{3.1}} in (\ref{text3_part4}). (It is only now that the speaker realizes the \isi{contrastive} potential between the negative effects of, respectively, the traffic and the cats on the well-being of the rats.) We, therefore, analyze the answer in (\ref{text3_part4}) as a \textsc{ct-f} structure.\footnote{But we should keep in mind that the assumption of a single focus spanning both the actor and the verb, hence a direct answer to question Q\textsubscript{1}, is also still an option.} This assumption is perhaps corroborated by a statement from one of the consultants who suggested the construction in A\textsubscript{{3.1}} and who explained her choice as follows: ``Both parts of the sentence [i.e.\ A\textsubscript{{3.1}} and A\textsubscript{4}] are about the cats, so we are talking about cats now", i.e.\ we have shifted the topic to the cats.

It is interesting to note that two people offered the construction given in (\ref{text3_part4}), but two other consultants chose a predicate-initial construction with \isi{undergoer voice} and the focal genitive-marked actor in the sentence-final position, as shown in (\ref{text3_part6}). 

\begin{exe}
\ex\label{text3_part6}\textbf{Alternative target sentence:}
\begin{xlist}
\exi{{>}{>}> Q\textsubscript{{1.1}}:}{\{What do cats do to rats?\}}
\end{xlist}
\begin{xlist}[{>}>{>}> A\textsubscript{{1.1}}:]
\exi{{>}>{>}> A\textsubscript{{1.1}}:}
\gll {{\ob}{\ob}h{\USSmaller}in{\USGreater}u-huli{\cb}\focus} rin {{\ob}ang} mga {daga{\cb}\topic} {{\ob}ng} mga {pusa{\cb}\ctopic{\cb}\sq}.\\
\textsc{\void{[[}<uv.rls>ipfv}-catch also \textsc{\void{[}nom} \textsc{pl} rat \textsc{\void{[}gen} \textsc{pl} cat\\
\glt `Also cats catch rats.'
\end{xlist}
\end{exe}

\noindent This construction, in which the focal genitive-marked actor is in the marked, sentence-final position and the topical \isi{undergoer} triggers \isi{undergoer voice}, is an example comparable to (\ref{martinez_b}) in \sectref{sect:is_tagalog}, in which the \textsc{ct} is clause-final. The speakers chose for the target sentence a structure identical to the sentence preceding it, i.e. they construed for the preceding sentence \textit{(the noise of the traffic made them nervous)} an \textsc{uv}-sentence with \textit{traffic} as genitive-marked \isi{actor argument}, so that the two affectors \textit{traffic} and \textit{cats} were contrasted. 

The focus sensitive particle \textit{rin} appears in both realizations, (\ref{text3_part4}) and (\ref{text3_part6}), right after the verb indicating that the verb is included in its scope. 
    
\subsection{Text 4 -- Intended: Actor (contrastive topic) Verb (background) Undergoer (focus)}\label{sect:ct_bg_f}

\begin{figure}
\framebox{
\footnotesize
\begin{tikzpicture}
\tikzset{inner xsep=0em,level distance=7.5em,sibling distance=0em,every tree node/.style={text width=6em,text centered}}
\Tree [.{Q\textsubscript{0}: \{What is the way things are?\}}
         [.{Q\textsubscript{{0.1}}: \{What is life in the wilderness like?\}}
            {A\textsubscript{{0.1}}: In the wilderness there are unfair laws.}
            [.{Q\textsubscript{1}: \{Who catches whom?\}}
              [.{Q\textsubscript{{1.1}}: \{Whom do lions catch?\}}
                {A\textsubscript{{1.1}}: Lions catch antelopes,}
              ]
              [.{Q\textsubscript{{1.2}}: \{Whom do sharks catch?\}}
                {A\textsubscript{{1.2}}: sharks catch tunafish.}
                [.{Q\textsubscript{{2}}: \{As for sharks and catching, what else is happening and with whom?\}}
                  [.{Q\textsubscript{{2.1}}: \{Whom are sharks getting caught by?\}}
                    {A\textsubscript{{2.1}}: They are furthermore getting caught by humans.}
                  ]
                ]
              ]
              [.{Q\textsubscript{{3}}: \{Where else can these laws of nature be observed?\}}
                {A\textsubscript{{3}}: Even in our cities this can be observed, too.}
              ]
              [.{Q\textsubscript{{1.3}}: \{Whom do cats catch?\}}
                   {\textbf{A\textsubscript{{1.3}}: Our cats catch rats.}}
              ]
            ]
      ]
     ]
\end{tikzpicture}
}
\caption{Discourse structure based on Text~4}
\label{fig:text4_analysis1}
\end{figure}

The fourth text (\figref{fig:text4_analysis1}) deals with life in the wilderness and different kinds of predators that chase and catch animals and are chased themselves.  Therefore, the text consists once again of a number of parallel sentences involving two variables \textsc{(ct+f)} that motivate the general question \textit{Who catches whom?} and the specific questions \textit{Whom do lions catch?} and \textit{Whom do sharks catch?} Note that these parallel \isi{contrastive} topic-focus sentences all show the same construction: \textit{ay}-inversion and \isi{actor voice}, as shown in (\ref{text4_part2}), in accordance with Hypothesis (i).

\begin{exe}
\ex\label{text4_part1}
\begin{xlist}
\exi{Q\textsubscript{0}:} \{What is the way things are?\}
\exi{> Q\textsubscript{{0.1}}:} \{What is life in the wilderness like?\}
\end{xlist}
\begin{xlist}[{>}> A\textsubscript{{0.1}}:]
\exi{{>}> A\textsubscript{{0.1}}:}
\gll {{\ob}{\ob}Sa} {kalikasan{\cb}\ctopic} ay {{\ob}may} mga hindi patas na {batas{\cb}\focus{\cb}\sq}.\\
\textsc{\void{[[}dat} wilderness \textsc{inv} {\void{[}exist} \textsc{pl} \text{neg} fair \textsc{lk} law\\ 
\glt `In the wilderness there are unfair laws.'
\end{xlist}
\end{exe}
\begin{exe}
\ex\label{text4_part2}
\begin{xlist}[{>}> A\textsubscript{{0.1}}:]
\exi{{>}> Q\textsubscript{1}:\hfill} \{Who catches whom?\}
\end{xlist}
\begin{xlist}[{>}{>}> Q\textsubscript{{1.2}}:]
\exi{{>}{>}> Q\textsubscript{{1.1}}:\hfill} \{Whom do lions catch?\}
\end{xlist}
\begin{xlist}[{>}>{>}> A\textsubscript{{1.2}}:]
\exi{{>}>{>}> A\textsubscript{{1.1}}:\hfill}
\gll {{\ob}{\ob}Ang} mga {leon{\cb}\ctopic} ay {nag-hu-huli}  {{\ob}ng} mga {bayawak{\cb}\focus{\cb}\sq}\\
\textsc{\void{[[}nom} \textsc{pl} lion \textsc{inv} \textsc{av.rls-ipfv}-catch \textsc{\void{[}gen} \textsc{pl} antelope\\
\glt `Lions catch antelopes,'
\end{xlist}
\begin{xlist}
\exi{{>}{>}> Q\textsubscript{{1.2}}:} \{Whom do sharks catch?\}
\end{xlist}
\begin{xlist}[{>}>{>}> A\textsubscript{{1.2}}:]
\exi{{>}>{>}> A\textsubscript{{1.2}}:}
\gll {{\ob}{\ob}Ang} mga {pating{\cb}\ctopic} ay nang-hu-huli {{\ob}ng} mga {tulingan{\cb}\focus{\cb}\sq}.\\
\textsc{\void{[[}nom} \textsc{pl} shark \textsc{inv} \textsc{av.rls-ipfv}-catch \textsc{\void{[}gen} \textsc{pl} tunafish\\
\glt `sharks catch tunafish.'
\end{xlist}
\end{exe}

\noindent The next sentence, A\textsubscript{{2.1}} in (\ref{text4_part3}), signals, both in its original formulation and in the \ili{Tagalog} version, a change in \isi{discourse} strategy: both the (nominative-marked) sharks and the activity of catching are kept up from A\textsubscript{{1.2}} in (\ref{text4_part2}) to A\textsubscript{{2.1}} in (\ref{text4_part3}). However, the sharks change their role from actor to \isi{undergoer}, which is expressed by a \isi{contrastive} change to \isi{undergoer voice}. Furthermore, a new actor \textit{(humans)} is introduced. We tentatively propose the -- unusual -- information structure shown in (\ref{text4_part3}), in which the \isi{voice} infix is assigned the function of a \isi{contrastive topic}, while the actor \textit{humans} is the focus. This also explains the slightly cumbersome formulation of Q\textsubscript{2} and Q\textsubscript{{2.1}}. Under this analysis, the \textit{ay}-inverted \isi{pronoun} is merely an \isi{aboutness topic}; a constellation for which we have not formulated any hypothesis. 

\begin{exe}
\ex\label{text4_part3}
\begin{xlist}[{>}>{>}> Q\textsubscript{2}:]
\exi{{>}>{>}> Q\textsubscript{2}:} \{As for sharks and catching, what else is happening and with whom?\}
\end{xlist}
\begin{xlist}
\exi{{>}>{>}{>}> Q\textsubscript{{2.1}}:} \{Whom are sharks getting caught by?\}
\end{xlist}
\begin{xlist}[{>}>{>}>{>}> A\textsubscript{{2.1}}:]
\exi{{>}>{>}>{>}> A\textsubscript{{2.1}}:}
\gll {{\ob}{\ob}Sila{\cb}\topic} ay {h{\ob}{\USSmaller}in{\USGreater}{\cb}\ctopic u-huli} din {{\ob}ng} mga {tao{\cb}\focus{\cb}\sq}.\\
\textsc{\void{[[}3pl.nom} \textsc{inv} \textsc{<uv>-ipfv}-catch also \textsc{\void{[}gen} \textsc{pl} man\\
\glt `They are furthermore getting caught  by humans'\\
Original: `and they happen to also get caught  by humans themselves.'
\end{xlist}
\end{exe}

\noindent The speaker then jumps back up in the tree with a comment about cities, (\ref{text4_part4}). What we see in A\textsubscript{{3}} is another \textit{ay}-inversion expressing a \isi{narrow contrastive focus} on a locative adjunct.\footnote{The verb \textit{observe} is treated as salient here, although it has not been mentioned explicitly in the previous \isi{discourse}.}

\largerpage[2]
\begin{exe}
\ex\label{text4_part4}
\begin{xlist}
\exi{{>}{>}> Q\textsubscript{{3}}:}{\{Where else can these laws of nature be observed?\}}
\end{xlist}
\begin{xlist}[{>}>{>}> A\textsubscript{{3}}:]
\exi{{>}>{>}> A\textsubscript{{3}}:\hfill}
\gll Kahit {{\ob}{\ob}sa} ating mga {lungsod{\cb}\focus} ay na-o-obersbahan din {{\ob}ito{\cb}\topic.{\cb}\sq}\\
Even \textsc{\void{[[}dat} our-\textsc{lk} \textsc{pl} city \textsc{inv} \textsc{uv-ipfv}-observe also \textsc{\void{[}dem.nom}\\
\glt `Even in our cities this can be observed, too.'\\
Original: `Even here in the city these cruel laws of nature can be observed.'
\end{xlist}
\end{exe}
 
\newpage  
\noindent The target sentence \textit{cats also catch rats} is then realized by three speakers as an \textit{ay}-inversion with \isi{actor voice}, i.e.\ the speaker is returning to the previously chosen \isi{discourse} strategy, expressed by Q\textsubscript{1}: \textit{\{Who catches whom?\}} and the follow-up subquestion Q\textsubscript{{1.3}} in (\ref{text4_part5}). 

\begin{exe}
\ex\label{text4_part5}\textbf{Target sentence:}
\begin{xlist}
\exi{{>}{>}> Q\textsubscript{{1.3}}:}{\{Whom do cats catch?\}} 
\end{xlist}
\begin{xlist}[{>}>{>}> A\textsubscript{{1.3}}:]
\exi{{>}>{>}> A\textsubscript{{1.3}}:}
\gll {{\ob}{\ob}Ang} nating mga {pusa{\cb}\ctopic} ay  nang-hu-huli {{\ob}ng} mga {daga{\cb}\focus{\cb}\sq}.\\ 
\textsc{\void{[[}nom} our \textsc{pl} cat \textsc{inv} \textsc{av.rls-ipfv}-catch \textsc{\void{[}gen} \textsc{pl} rat\\
\glt `Our cats catch rats.'
\end{xlist}
\end{exe}

\noindent The fourth speaker suggested the \textit{ang}-inversion in (\ref{text4_part6}) -- recall example (\ref{nom_inv1}) in \sectref{sect:is_tagalog} -- thus either expressing the thought that the only ones who catch rats are cats -- as indicated below -- or, alternatively, that the only act of catching that takes place in the city is between cats and rats. The first option would mean that, in violation of the \textsc{Q-Givenness} principle (\sectref{sect:principles}), the rats would have to be accommodated as given. The second solution would inevitably mean that the \textit{ang}-inversion construction is not restricted to exhaustive narrow \isi{argument focus}, but may also be used for (exhaustive) complex focus spanning both the actor and the \isi{undergoer}. The data here are too limited to finally settle this question. However, given that we are not aware of languages in which cleft sentences are restricted to narrow argument marking, the second solution seems quite plausible.

\begin{exe}
\ex\label{text4_part6}\textbf{Alternative target sentence:}
\begin{xlist}
\exi{{>}{>}> Q\textsubscript{4}:} \{Who catches rats (in the city)?\}
\end{xlist}
\begin{xlist}[{>}>{>}> A\textsubscript{4}:]
\exi{{>}>{>}> A\textsubscript{4}:} 
\gll {{\ob}{\ob}Ang} nating mga {pusa{\cb}\focus} ang  {nang-hu-huli} {{\ob}ng} mga {daga{\cb}\topic{\cb}\sq}.\\
\textsc{\void{[[}nom} our \textsc{pl} cat \textsc{nom} \textsc{<av.rls>ipfv}-catch \textsc{\void{[}gen} \textsc{pl} rat\\ 
\glt `(Here) it is our cats that catch rats.'\\
Original sentence: `(Our domestic) cats also chase and catch rats.' 
\end{xlist}
\end{exe}


\subsection{Text 5 -- Intended: Sentence focus, i.e.\ Actor (focal) Verb (focal) Undergoer (focal); Partly realized: Actor (contrastive topic) Verb (focal) Undergoer (focal)}\label{sect:text_new}

\begin{figure}
\framebox{
\small
\begin{tikzpicture}
\tikzset{inner xsep=0em,level distance=5em,sibling distance=0em,every tree node/.style={text width=8em,text centered}}
\Tree [.{Q\textsubscript{0}: \{What is the world like?\}}
         [.{Q\textsubscript{{0.1}}: \{What do I see?\}}
            {A\textsubscript{{0.1}}: When I look out of the window, I see only unhappiness and turmoil.}
            [.{Q\textsubscript{1}: \{Who does what?\}}
              [.{Q\textsubscript{{1.1}}: \{What do the neighbours do?\}}
                {A\textsubscript{{1.1}}: The neighbours fight with the children.}
              ]
              [.{Q\textsubscript{{1.2}}: \{What do the bears do?\}}
                {A\textsubscript{{1.2}}: Bears persecute birds}
              ]
              [.{Q\textsubscript{{1.3}}: \{What do the cats do?\}}
                {A\textsubscript{{1.3}}: and cats catch rats.}
              ]
             ]
         ]
      ]
\end{tikzpicture}
}
\caption{Discourse structure based on Text~5}
\label{fig:text5_analysis1}
\end{figure}

\noindent Finally, we tried to construe a context for an all-new sentence, in order to elicit a clause-focus construction. In the text shown in \figref{fig:text5_analysis1}, a number of different scenes witnessed from a window were listed. Every sentence contained a new actor, a new verb and a new \isi{undergoer}. 

In the first \isi{assertion} A\textsubscript{{0.1}} in (\ref{text5_part1}), the \textit{ang}-inverted narrow (exhaustive) focus on \textit{panay-away at kaguluhan} `constant fighting and turmoil' is motivated by the focus-sensitive particle \textit{only}, i.e.\ the question Q\textsubscript{{0.1}}: \textit{\{What do I see?\}} is accommodated. Moreover, the \textit{when}-\isi{clause} is interpreted as a relevance condition which is not at issue. 

\begin{exe}
\ex\label{text5_part1}
\begin{xlist}
\exi{Q\textsubscript{0}:} \{What is the world like?\}
\exi{> Q\textsubscript{{0.1}}:} \{What do I see?\}
\end{xlist}
\begin{xlist}[{>}> A\textsubscript{{0.1}}:]
\exi{{>}> A\textsubscript{{0.1}}:}
\gll {{\ob}Pag} t{\USSmaller}um{\USGreater}i-tingin ako sa {bintana,{\cb}\nai} {{\ob}{\ob}panay-away} at    {kaguluhan{\cb}\focus} lang    ang    na-ki-kita {ko{\cb}\sq}.\\
{\void{[}when} \textsc{<av.rls>ipfv}-look \textsc{1sg.nom} \textsc{dat} window {\void{[[}constant fight} and turmoil only  \textsc{nom} \textsc{uv-ipfv}-see \textsc{1sg.gen}\\
\glt `When I look out of the window, I see only constant fighting and turmoil.'
\end{xlist}
\end{exe}

\noindent This sentence is followed by yet another  sequence of \textit{ay}-inversions (i.e.\ \textsc{ct-f} structures), describing who \textsc{(ct)} is performing which act of violence:  the focused elements are, therefore, the \textsc{av}-marked predicate and the \isi{undergoer}. The construction fits the pattern described in Hypothesis (i).\footnote{What this example shows is that the construction of a text that consists of only new and unrelated sentences is in fact a very unnatural task. Instead, human interpreters will seize every opportunity to bring structure (here: \textsc{ct-f} pairs) -- and therefore coherence -- into what was originally intended to be an unstructured sequence of sentences.}

\newpage 
\begin{exe}
\ex\label{text5_part2a}
\begin{xlist}
\exi{{>}> Q\textsubscript{1}:} \{Who does what?\}
\end{xlist}
\begin{xlist}[{>}{>}> Q\textsubscript{{1.3}}:]
\exi{{>}{>}> Q\textsubscript{{1.1}}:\hfill} \{What do the neighbours do?\}
\end{xlist}
\begin{xlist}[{>}>{>}> A\textsubscript{{1.3}}:]
\exi{{>}>{>}> A\textsubscript{{1.1}}:\hfill}
\gll {{\ob}{\ob}Ang} mga {kapitbahay{\cb}\ctopic} ay {{\ob}nag-a-away} ng mga {bata{\cb}\focus{\cb}\sq}.\\
\textsc{\void{[[}nom} \textsc{pl} neighbour \textsc{inv} \textsc{\void{[}av.rls-ipfv}-fight \textsc{gen} \textsc{pl} child\\
\glt `The neighbours fight with the children.' 
\end{xlist}
\begin{xlist}[{>}{>}> Q\textsubscript{{1.3}}:]
\exi{{>}{>}> Q\textsubscript{{1.2}}:\hfill} \{What do the bears do?\}
\end{xlist}
\begin{xlist}[{>}>{>}> A\textsubscript{{1.3}}:]
\exi{{>}>{>}> A\textsubscript{{1.2}}:\hfill}
\gll {{\ob}{\ob}ang} mga {oso{\cb}\ctopic} ay {{\ob}t{\USSmaller}um{\USGreater}u-tusig} ng mga {birds{\cb}\focus{\cb}\sq}\\
 \textsc{\void{[[}nom} \textsc{pl} bear \textsc{inv} \textsc{\void{[}<av.rls>ipfv}-persecute  \textsc{gen} \textsc{pl} bird\\
\glt `Bears persecute birds' 
\end{xlist}
\end{exe}

\begin{exe}
\ex\label{text5_part2}\textbf{Target sentence:}
\begin{xlist}
\exi{{>}{>}> Q\textsubscript{{1.3}}:} \{What do the cats do?\}
\end{xlist}
\begin{xlist}[{>}>{>}> A\textsubscript{{1.3}}:]
\exi{{>}>{>}> A\textsubscript{{1.3}}:}
\gll at {{\ob}{\ob}ang} mga {pusa{\cb}\ctopic} ay {{\ob}nang-hu-huli} ng mga {daga{\cb}\focus{\cb}}.\\
 and \textsc{\void{[[}nom} \textsc{pl} cat \textsc{inv} \textsc{\void{[}av.rls-ipfv}-catch \textsc{gen} \textsc{pl} rat\\
\glt `and cats catch rats.'
\end{xlist}
\end{exe}

\noindent Three speakers continued with a third \textit{ay}-inverted \textsc{ct-f} \isi{assertion} for the target sentence, as shown A\textsubscript{{1.3}} of (\ref{text5_part2}). Only one speaker chose an all-focus existential construction for the target sentence. 

\begin{exe}
\ex\label{text5_part3}\textbf{Alternative target sentence:}
\begin{xlist}
\exi{{>}> Q\textsubscript{2}:} \{What other violent things are happening?\}
\end{xlist}
\begin{xlist}[{>}{>}> A\textsubscript{2}:]
\exi{{>}{>}> A\textsubscript{2}:}
\gll {{\ob}{\ob}May} mga pusang h{\USSmaller}um{\USGreater}a-habol at p{\USSmaller}um{\USGreater}a-patay ng mga inosente-ng mga {daga{\cb}\focus{\cb}\sq}.\\
{\void{[[}exist} \textsc{pl} cat-\textsc{lk} \textsc{<av>-ipfv}-catch and \textsc{<av>-ipfv}-kill \textsc{gen} \textsc{pl} innocent-\textsc{lk} \textsc{pl} rat\\
\glt `There are cats chasing and killing innocent rats.'\\
\end{xlist}
\end{exe}

\section{Summary of findings, conclusions}

\begin{table}
\caption{Information structure and morphosyntactic choices (voice, constituent order and inversion): target realizations in five sample texts}
\label{tab:overview1}
\begin{tabularx}{\textwidth}{l>{\scshape}CCCc} 
\lspbottomrule
Structure & {\upshape Voice} & Example & Hypothesis & Text \\ 
\midrule 
{[V UG]}\focus\ ACT\topic  & av & (\ref{text1_part3}) & (iii) &1 \\
\tablevspace
ACT\focus\ ay V\bg\ UG\topic & uv & (\ref{text2_part3}) & (v) & \multirow{2}{*}{2} \\ 
ACT\ctopic\ ay V\focus\ UG\topic & av & (\ref{text2_part2}) & (vi) & \\ 
\tablevspace
ACT\ctopic\ ay V\focus\ UG\topic & av & (\ref{text3_part4}) & (vi) & \multirow{2}{*}{3} \\ 
V\focus\ UG\topic\ ACT\ctopic & uv & (\ref{text3_part6}) & -- &  \\ 
\tablevspace
ACT\ctopic\ ay V\bg\ UG\focus & av & (\ref{text4_part5}) & (i) & \multirow{2}{*}{4} \\
ACT\focus\ ang V\bg\ UG\subs{T/F(?)} & av & (\ref{text4_part6}) & -- & \\ 
\tablevspace
ACT\ctopic\ ay [V UG]\focus & av & (\ref{text5_part2}) & (i) & \multirow{2}{*}{5} \\
{[may ACT V UG]}\focus & av & (\ref{text5_part3}) & (i) & \\
\lspbottomrule
\end{tabularx}
\end{table}

\noindent \tabref{tab:overview1} sums up our findings with respect to the target-sentence translations discussed in the paper. Note that the other (non-target) sentences discussed in the previous sections are not mentioned in the table, although their analyses, too, are in accordance with the described hypotheses. 

In four of the texts (except for Text~1) the actor was focal (which includes \isi{contrastive} topics) and thus IS-prominent, in two of them (Text~2 and 3) the \isi{undergoer} was topical/\isi{backgrounded} and thus IS-prominent.  In our examples, \isi{actor voice} was the preferred choice. As \tabref{tab:overview1} shows, with the exception of Text 1, the \isi{actor argument} was always conceived of as a \isi{contrastive topic} or a focus. Contrastive topics (and, otherwise, narrow foci) were typically 
\textit{ay}-inverted (Texts~2--5). Notable exceptions were the extraposed \textsc{ct} in (\ref{text3_part6}), the existential construction in (\ref{text5_part3}) as well as the \textit{ang}-inversion in (\ref{text4_part6}), whose information-structural analysis remains uncertain. \isi{Undergoer} \isi{topicality} did not generally lead to the selection of \isi{undergoer voice} in the presence of a focal actor (Texts~2 and 3). The current data lead us, thus, to the conclusion that focality of the actor is definitely more salient, and has a greater morphosyntactic effect, than \isi{topicality} of the \isi{undergoer}.

In general, it has become clear that in order to describe the information-structural impact on \isi{voice selection}, a more comprehensive approach is necessary rather than simply considering \isi{givenness} and newness of arguments. This paper\footnote{See also the paper by \citetv{riear18c}.} is the first to apply the new \textit{QUD-tree} method described in \cite{riear18b} to \ili{Austronesian} language data. It, therefore, demonstrates a completely new way of studying the information structure of a lesser-described language on the basis of textual corpus data.

With respect to our hypotheses, we can specifically state the following based on our case study: Hypothesis (i) was confirmed in the data.

\hypo{i}{Actor$_{\text{\footnotesize F/CT}}$ \wupperhand ay \isi{Undergoer}\focus\ $\rightarrow$ Actor Voice}

\noindent However, we must add that, in our data, basic sentences structure was not chosen at all 
to encode two focused (or \textsc{ct}) core arguments. Rather we found the \textit{ay}-inversion construction to be the predominant pattern (Texts~4 and 5).

  
\hypo{ii}{Actor\topic\ \isi{Undergoer}\topic\ $\rightarrow$ \isi{Undergoer} Voice}

\noindent Hypothesis (ii) was not truly investigated here, as basically all \isi{undergoer voice} sentences found and discussed in the literature are of this type. Therefore, it was not the most interesting case to look at. Our target sentence did not appear in a narrow verb-focus or adjunct focus context, and apart from one sentence with a focal locative adjunct, (\ref{text4_part4}) in Text~4, which indeed showed \isi{undergoer voice}, we have nothing to add to this particular issue. 

\hypo{iii}{Actor\topic\ \isi{Undergoer}\focus\ $\rightarrow$ decided at different level}

\noindent Since Hypothesis (iii) left it open whether a default information structure would lead to actor or \isi{undergoer voice}, it was not actually challenged by our data. In the undergoer-verb-focus scenario (Text~1), the participants provided actor-\isi{voice} sentences; more specifically, they chose 
the expected unmarked verb-initial \isi{word order}.

\hypo{iv}{\isi{Undergoer}\topic\ Actor\focus\ (default order)  $\rightarrow$ Actor Voice}

\noindent Hypothesis (iv) was discussed in \sectref{sect:is_tagalog} and eventually restricted to cases with default \isi{constituent order}, thus excluding patterns with \textit{ay-}inversion and extraposed actors. Remember that \isi{narrow actor focus} in our data was expressed by means of \textit{ay-}inversion. Hence, Hypothesis (iv) did not apply to any of the cases found in the data.

Based on the great variety of cases involving \textit{ay}-inversion in our data, we formulated two more hypotheses regarding \textit{ay}-inverted actors and clause-final topical undergoers.

\hypo{v}{Actor\focus\ +ay Verb\bg\ \isi{Undergoer}\topic\ $\rightarrow$ \isi{Undergoer} Voice}

\noindent Hypothesis (v) specifies one way of realizing \isi{narrow actor focus}, hence, the situation intended in Text~2, which was found in sentence (\ref{text2_part3}).

\hypo{vi}{Actor\ctopic\ +ay Verb\focus\ \isi{Undergoer}\topic\ $\rightarrow$ Actor Voice} 

\noindent Finally, Hypothesis (vi) could be confirmed for all instances exhibiting the respective syntactic-pragmatic pattern, which, however, were produced by the consultants more or less by accident, since the \textsc{ct-f} structures they came up with deviated from the originally intended in\-for\-ma\-tion-struc\-tu\-ral constellations. Notably, if we examine all examples that exhibit an \textit{ay}-inverted \isi{contrastive topic} and a focus on either the verb, the \isi{undergoer} or both, we always witness \isi{actor voice}. 

The lesson we learn from our approach is that, on the one hand, it is quite difficult to steer participants to produce a very specific information-structural pattern and obtain a natural \isi{discourse} at the same time. On the other hand, the QUD-tree method allows us to fruitfully analyze and interpret the actually produced data despite the deviations we gained.

In the current study we have presented many new facts about information structure marking in \ili{Tagalog} and, in particular, about its relation to \isi{voice} and inversion. We were able to specify a number of detailed hypotheses, and data that match them. We must leave it to future work, though, to test these hypotheses on a bigger scale and to identify a more general explanation why certain \isi{voice} forms were chosen under specific information-structural constellations.

\section*{Acknowledgements}
The authors would like to thank the editors of this volume and two anonymous reviewers. We also owe many thanks to the LingDy project in Tokyo for invitations to several workshops on the information structure of \ili{Austronesian} languages, held at the Research Institute for Languages and Cultures of Asia and Africa (ILCAA) at Tokyo University of Foreign Studies. Furthermore, financial support of Deutsche Forschungsgemeinschaft (DFG) is kindly acknowledged. The authors are funded via the Stuttgart Sonderforschungsbereich 732, Project A6 and the D\"usseldorf Sonderforschungsbereich 991, Project D4. Special thanks go to R.D.\ Van Valin, Jr., Katalin Balogh and Patrick Nuhn for comments on earlier versions of the paper.

 
\section*{Abbreviations}

\begin{multicols}{2}
	\begin{tabbing}
		glossgloss \= \kill
		\textsc{\sq} \> \isi{focus domain}\\
		\textsc{act} \> actor\\
		\textsc{appl$_\text{instr}$} \> instrumental applicative\\
		\textsc{av} \> \isi{actor voice}\\
		\textsc{bg} \> background\\
		\textsc{bv} \> beneficiary \isi{voice}\\
		\textsc{com} \> comitative\\
		\textsc{ct} \> \isi{contrastive topic}\\
		\textsc{dat} \> dative\\
		\textsc{exist} \> existential\\
		\textsc{f} \> focus\\
		\textsc{gen} \> genitive\\
		\textsc{gv} \> goal \isi{voice}\\
		\textsc{inv} \> inversion particle\\
		\textsc{ipfv} \> imperfective\\
		\textsc{iv} \> instrumental \isi{voice}\\
		\textsc{lk} \> linker\\
		\textsc{neg} \> negation\\
		\textsc{nom} \> \isi{nominative}\\
		\textsc{obl} \> oblique\\
		\textsc{pl} \> plural\\
		\textsc{q} \> question marker\\
		\textsc{rls} \> realis\\
		\textsc{sg} \> singular\\
		\textsc{t} \> \isi{aboutness topic}\\
		\textsc{tv} \> theme \isi{voice}\\
		\textsc{ug} \> \isi{undergoer}\\
		\textsc{uv} \> \isi{undergoer voice}
	\end{tabbing}
\end{multicols}

\sloppy
\printbibliography[heading=subbibliography,notkeyword=this]

\end{document}
